% ===========================================================================
%  IntroAct-TS —— ICLR 2027 投稿正文
% ---------------------------------------------------------------------------
%  版本      : v4.4（2026-09-17 首版）
%  模板      : 用户提供的 ICLR 2027 style files（未作任何改动）
%  上一版    : 无（v4.3.1-r5 只有 Markdown 草稿，本文件是第一个 LaTeX 正文）
%
%  重要约定
%  ---------
%  1. 所有尚未在服务器上跑出的数字，一律写成占位符 \ph{KEY}，渲染为灰色
%     [KEY]。绝不允许用估计值、旧协议值或"看起来合理"的数填充。
%     占位符命名规范见 latex/README.md。
%  2. 只有 v4.3.1-r5 台账里已经独立审核过的历史数字，才允许以真实数值出现在
%     Appendix G，并且必须显式标注其属于旧协议、不得与 v4.4 主表比较。
%  3. 正文只保留 4.1-4.6 六个小节，效率并入 4.2 末尾。其余全部进 Appendix。
%  4. 结论措辞受 docs/v431_r5_* 与规划 §41 的 claim 边界约束：在独立主实验完成
%     前不写 SOTA、不写 significantly outperforms、不写 safe/robust/generalizes。
% ===========================================================================
\documentclass{article}

\usepackage{iclr2027_conference,times}
\input{math_commands.tex}

\usepackage{hyperref}
\usepackage{url}
\usepackage{graphicx}
\usepackage{booktabs}
\usepackage{amsmath}
\usepackage{amssymb}
\usepackage{multirow}
\usepackage{array}
\usepackage[table]{xcolor}

% ---------------------------------------------------------------- 版式与宏
\definecolor{phgray}{gray}{0.45}
\definecolor{bestgray}{gray}{0.94}

% 未运行结果的占位符。KEY 一律用下划线分层命名（如 MAIN_OVERALL_MASE），
% 下划线是数学下标字符，因此必须 \detokenize 成普通字符后才能进 \textsf。
% 填充时只替换 KEY，不要改 \ph 本身。
\newcommand{\ph}[1]{\textcolor{phgray}{\textsf{[\detokenize{#1}]}}}
% 表格中强调：第一名加粗，第二名下划线，Oracle 用斜体灰字。
\newcommand{\best}[1]{\textbf{#1}}
\newcommand{\second}[1]{\underline{#1}}
\newcommand{\oracle}[1]{\textit{\textcolor{phgray}{#1}}}
\newcommand{\introact}{\textsc{IntroAct-TS}}
\newcommand{\tsfm}{TSFM}
\newcommand{\calA}{\mathcal{A}}
\newcommand{\calB}{\mathcal{B}}
\newcommand{\calN}{\mathcal{N}}
\newcommand{\loss}{\ell}

\newcommand{\fix}{\marginpar{FIX}}
\newcommand{\new}{\marginpar{NEW}}

% ICLR 正文硬上限 9 页，浮动体间距默认偏松，这里收紧以腾出正文空间。
\setlength{\textfloatsep}{7pt plus 2pt minus 2pt}
\setlength{\floatsep}{7pt plus 2pt minus 2pt}
\setlength{\intextsep}{7pt plus 2pt minus 2pt}
% 正文有 6 张表 + 5 张图，默认每页浮动体上限会让它们堆积到最后一页。
% 放宽每页浮动体配额，让它们均匀分布在第 1-9 页内。
\setcounter{topnumber}{4}
\setcounter{bottomnumber}{3}
\setcounter{totalnumber}{5}
\renewcommand{\topfraction}{0.92}
\renewcommand{\bottomfraction}{0.65}
\renewcommand{\textfraction}{0.06}
\renewcommand{\floatpagefraction}{0.72}

\title{\introact{}: Counterfactual Replay for Selective Data Governance \\
in Time-Series Foundation Models}

\author{Anonymous Authors\\Paper under double-blind review}

%\iclrfinalcopy
\begin{document}

\maketitle

\begin{abstract}
Time-series foundation models are increasingly deployed without task-specific parameter
updates, yet their forecasting contexts may contain missing variables, delayed observations,
or test-time missing blocks. Existing incomplete-input forecasting methods largely focus on
reconstructing unavailable information or on training architectures that remain robust to
missingness. We study a complementary question for frozen forecasters: given several legal
ways of handling an incomplete context, which intervention should be applied to the current
forecasting episode, if any? We introduce \introact{}, a selective data-governance framework
that replays historical pseudo-deployments to obtain the realised forecasting consequences of
candidate interventions, and derives its decision rule from evidence aligned with intervention
ranking rather than with reconstruction fidelity. \textsc{Keep} is treated as a first-class
action, so the method leaves an input unchanged whenever the available evidence does not
support a switch. \introact{} never fine-tunes the frozen backbone, never ensembles forecast
outputs, and returns a forecast produced from exactly one executed input version. Across
\ph{N_DATASETS} datasets, four missingness patterns, and two forecasting horizons, with two
frozen backbones used for development and a third family held out from method design,
\introact{} reaches \ph{MAIN_OVERALL_MASE} source-macro MASE against \ph{STRONG_BASELINE} at an
intervention rate of \ph{MAIN_INTERVENTION_RATE}, while the harmful-intervention diagnostics in
\S\ref{sec:exp-opportunity} remain \ph{NOOP_HARM}. These results suggest that incomplete-input
handling is better posed as a per-episode decision about when \emph{not} to act than as a
mandatory preprocessing step. \ph{ABSTRACT_CLOSING}.
\end{abstract}

% ===========================================================================
\section{Introduction}
\label{sec:intro}

Time-series foundation models (\tsfm{}s) now produce competitive zero-shot forecasts without
task-specific training~\citep{ansari2024chronos,chronosbolt2025,das2024timesfm,woo2024moirai,
chronos2_2025}, which makes them attractive as fixed components of deployed forecasting
systems. In deployment, however, the context handed to such a model is frequently incomplete.
A sensor drops out for a stretch of time; a reporting channel delivers a variable several
steps late; variables are sampled asynchronously and a few have not arrived at the forecast
origin; only part of an acquisition batch is available. The resulting input is not
\emph{wrong} --- it is \emph{partial}. An observation that has already arrived and is valid
must not be overwritten, and an incomplete context is a structural property of acquisition
rather than a defect to be scrubbed away.

The classical response to incompleteness is imputation. Reconstruction-oriented methods
recover missing entries as accurately as possible --- BRITS~\citep{cao2018brits} through
bidirectional recurrent dynamics, CSDI~\citep{tashiro2021csdi} through conditional
score-based diffusion, and generative or self-attention variants in the same
spirit~\citep{yoon2018gain,du2023saits} --- and are scored on reconstruction error. A more
recent line argues that this objective is a proxy. Task-oriented imputation
evaluation~\citep{wang2024taskoriented} estimates how much a given imputation strategy
improves a downstream task and shows that strategy rankings under reconstruction and under
task utility can disagree. TOI-VSF~\citep{hao2025toivsf} makes the point concrete for
\emph{variable subset forecasting} (VSF), where whole variables are absent at inference
time, by training the imputation module to support the forecasting objective directly rather
than to recover the missing series precisely. The question this line answers is
\emph{how should incomplete variables or observations be reconstructed or modelled so that
forecasting remains accurate}.

That question has since been pursued directly, and the resulting methods are progressively
more task- and deployment-aware. GIMCC~\citep{hao2025gimcc} couples generative imputation
with multi-level causal consistency so that filled-in variables preserve the causal structure
the forecaster relies on; SRDI~\citep{xu2026srdi} targets the distribution shift induced by
variable absence and imputes diffusively in a shift-resilient space;
ChannelTokenFormer~\citep{jang2026channeltokenformer} instead redesigns the forecasting
architecture so that channel dependency, sampling asynchrony, and test-time missing blocks
are addressed jointly rather than in isolation. Against generic imputation, these methods are
already aimed at incomplete-input forecasting, so \emph{task-oriented handling of missing
data} is no longer a contribution by itself. VIDA~\citep{liang2025vida} is a closely related
VSF framework that reframes the problem as cross-domain knowledge transfer; we discuss it
alongside the other extended comparisons in Appendix~\ref{app:extended}.

We study a question this line leaves open. The methods above answer how incomplete
observations should be reconstructed or modelled, they are evaluated by a score aggregated
over a corpus, and what they return is a single repaired or robustly modelled series. A
deployed system must instead decide, for \emph{this} episode, whether to intervene at all
and, if so, which of several legal ways of handling the incomplete input to use. Three
properties make that decision non-trivial. \textbf{(i)}~\emph{Intervention utility is
heterogeneous}: the same repair lowers forecasting loss on one episode and raises it on
another, so a strategy-level average does not license an action on a single window.
\textbf{(ii)}~\emph{No single global repair policy suffices}, because the identity of the
best action varies with the episode, the dataset, and the frozen backbone.
\textbf{(iii)}~\emph{A wrong intervention can be worse than doing nothing}, so unconditional
repair is not merely suboptimal but actively harmful on a non-trivial fraction of requests.
Reconstruction-oriented objectives and task-wide adaptation both optimise an aggregate, and
neither resolves per-episode action selection.

We therefore formalise incomplete-input handling for frozen \tsfm{}s as \emph{selective data
governance} and propose \introact{}, organised as three stages. \emph{Counterfactual Replay}
builds pseudo-deployments from historical windows whose futures are already legal history: it
injects the deployment missingness mechanism, actually executes every legal action with the
frozen forecaster, and records the realised forecasting loss of each action.
\emph{Regret-Aware Intervention Scoring} turns those records into evidence aligned with the
discrete decision we have to make, rather than into a point-wise regression of every action's
gain. \emph{Selective Intervention} treats \textsc{Keep} as an ordinary action, so abstention
is a member of the decision space rather than a separate safety layer. The returned forecast
always comes from exactly one actually executed input version; the backbone is never
fine-tuned, never ensembled, and never corrected.

\paragraph{Contributions.}
\textbf{(1)~Selective data governance.} We formalise incomplete-input handling for frozen
\tsfm{}s as a per-episode decision between executing one legal intervention and leaving the
input unchanged, and show that a feasible intervention is not always a beneficial one
(\S\ref{sec:exp-opportunity}). \textbf{(2)~Counterfactual replay.} We build a replay bank from
TRAIN windows by re-enacting the deployment missingness mechanism and executing the legal
action catalog with the real frozen forecaster, so that the stored supervision is a realised
task consequence rather than a reconstruction score (\S\ref{sec:method-replay},
\S\ref{sec:exp-main}). \textbf{(3)~Regret-aware selective decision.} We derive the decision
from historical action regret and pairwise preference instead of from a point-wise gain
regression, which lets \textsc{Keep} compete on equal terms and makes abstention a natural
outcome rather than an imposed guard (\S\ref{sec:method-scoring}, \S\ref{sec:method-decision},
\S\ref{sec:exp-robust}).

\begin{figure}[t]
  \centering
  \includegraphics[width=0.80\textwidth]{figure/fig1_selective_governance.pdf}
  \caption{Intervention utility is episode dependent. \textbf{(a)}~Three incomplete-context
  episodes with the same missingness mechanism and the same frozen forecaster, but different
  best actions: on episode A the reference input is already best, on B a different legal
  action wins, and on C yet another one does. \textbf{(b)}~Historical TRAIN windows are
  replayed as pseudo-deployments: the deployment missingness mechanism is injected, every
  legal action is actually executed with the frozen forecaster, and the realised action losses
  are recorded. \textbf{(c)}~A new request is scored against that replay evidence, and the
  system either executes the best-ranked legal action or returns the untouched reference
  forecast.}
  \label{fig:concept}
\end{figure}

% ===========================================================================
\section{Related Work}
\label{sec:related}

\subsection{Reconstruction-Oriented Missing Data Modeling}
Imputation is a mature line. BRITS~\citep{cao2018brits} performs bidirectional recurrent
imputation and was an early demonstration that temporal consistency improves recovery;
CSDI~\citep{tashiro2021csdi} replaces the deterministic mapping with a conditional
score-based diffusion model; and GAIN~\citep{yoon2018gain} and SAITS~\citep{du2023saits}
explore adversarial and self-attention variants of the same objective. All of these methods
are scored on reconstruction error against the artificially hidden values. We keep this line
as background only: our object of study is not a better reconstruction model, and no number
in this paper rewards reconstruction quality for its own sake.

\subsection{Task-Oriented Incomplete-Input Forecasting}
\label{sec:related-taskoriented}
A second line moves the objective from recovery to downstream utility, and its methods have
become progressively closer to our setting. Task-oriented imputation
evaluation~\citep{wang2024taskoriented} estimates how much a given imputation strategy
improves a downstream task without retraining the task model, and reports that strategy
rankings under reconstruction and under task utility can disagree. TOI-VSF~\citep{hao2025toivsf}
carries this into \emph{variable subset forecasting}, where entire variables are absent at
inference: it pairs a forecasting-model-agnostic self-supervised imputation module with joint
imputation-and-forecasting learning, on the argument that precise recovery is not the
objective. GIMCC~\citep{hao2025gimcc} adds multi-level causal consistency between the
generated values and the forecasting target, so that imputation preserves the structure the
forecaster depends on. VIDA~\citep{liang2025vida} reframes VSF as cross-domain knowledge
transfer and adapts representations from the complete-variable domain.
SRDI~\citep{xu2026srdi} addresses the distribution shift that variable absence itself induces,
imputing diffusively in a shift-resilient space. The trajectory is therefore: task-oriented
evaluation, then variable-subset forecasting, then task-driven reconstruction, then
shift-aware VSF.

What this line still produces is a \emph{strategy-level} decision and a single repaired or
robustly modelled series. Our question is finer-grained: given several legal ways of handling
the current incomplete input, should the system intervene at all, and which intervention
serves \emph{this} frozen forecaster on \emph{this} episode. We also constrain the returned
forecast to come from one actually executed input version, so no output fusion or post-hoc
correction can contribute to a reported number.

\subsection{Robust Forecasting with Missing and Asynchronous Inputs}
\label{sec:related-robust}
A third line changes the forecasting model rather than the data.
ChannelTokenFormer~\citep{jang2026channeltokenformer} is a Transformer framework designed so
that inter-channel dependency, channel-wise asynchronous sampling, and test-time missing
blocks are handled jointly instead of in isolation, and it is evaluated on both public
benchmarks and a private industrial dataset. Methods in this family are complete forecasting
architectures: they are trained for the task and their parameters are part of the contribution.

\introact{} differs from all three lines on a single axis. It does not propose a new
forecasting backbone, does not require training or modifying the frozen \tsfm{}, and does not
output a fused or corrected series. It instead operates inside a fixed catalog of legal
interventions and makes an \emph{episode-level selection} among them, with the untouched
input available as one of the choices. Because an arm's value here is the realised loss of a
frozen, non-differentiable forecaster, the decision layer also differs from conservative
policy selection in contextual bandits and from off-policy evaluation: there is no designed
reward, no simulator, and no policy parameters to fit. We develop that contrast in
Appendix~\ref{app:ope-positioning} and keep it out of the main text because the selector we
report is not a policy-learning method.

% ===========================================================================
\section{Problem Formulation and Method}
\label{sec:method}

\begin{figure}[t]
  \centering
  \includegraphics[width=0.62\textwidth]{figure/fig2_introact_architecture.pdf}
  \caption{Architecture of \introact{}. \textbf{Offline (top)}: complete historical TRAIN
  windows are turned into controlled incomplete contexts by injecting the deployment
  missingness mechanism; every legal intervention is executed with the frozen forecaster and
  the realised action losses populate a counterfactual replay bank. The offline path is the
  only place where a future is read. \textbf{Online (bottom)}: a deployment request exposes an
  incomplete context only, from which an observable state is computed; that state is scored
  against the frozen bank, and one legal action --- possibly \textsc{Keep} --- is executed with
  the same frozen forecaster. Neither a future target nor a non-selected candidate's forecast
  is ever accessed online.}
  \label{fig:arch}
\end{figure}

\subsection{Problem Formulation}
\label{sec:method-problem}
Let $X \in \mathbb{R}^{C \times L}$ be the multivariate context available at the forecast
origin, with observation mask $M \in \{0,1\}^{C \times L}$, and let
$Y \in \mathbb{R}^{C \times H}$ be the future target, which is \emph{never} available to
the decision procedure. Let $F$ be a frozen time-series foundation model with fixed
parameters, mapping an input to a forecast $p = F(X)$. Let
$\calA = \{a_0, a_1, \dots, a_K\}$ be a fixed catalog of legal governance actions, where
each $a$ materialises a legal input version $X^{a}$ from $(X, M)$ and $a_0$ is the
\emph{reference} action (the untouched input, i.e.\ \textsc{Keep}). Let $\loss$ be a
forecasting loss. The task is to choose
\begin{equation}
  a^\star(X, M) \;\in\; \calA,
  \qquad
  a^\star \text{ must not depend on } Y ,
  \label{eq:objective}
\end{equation}
so as to improve $\loss(F(X^{a^\star}), Y)$ relative to the reference
$\loss(F(X^{a_0}), Y)$, while keeping the number of frozen-model calls small and avoiding
harmful interventions.

Three structural constraints distinguish this from ordinary imputation. First, \emph{no
valid observation may be overwritten}: every $X^{a}$ differs from $X^{a_0}$ only at
positions where $M = 0$. Second, \emph{the future is unavailable}: no candidate forecast
may be inspected before an action is chosen. Third, \emph{the untouched input is a
candidate like any other}: \textsc{Keep} belongs to $\calA$ and the decision rule of
\S\ref{sec:method-decision} can select it, so leaving the input unchanged is an outcome of
the same ranking rather than an exception path bolted onto it.

\subsection{Counterfactual Replay}
\label{sec:method-replay}
\paragraph{Terminology.}
We use \emph{counterfactual} in the offline-evaluation sense, not the causal-identification
sense: for $a \neq a_0$, $L_{i,a}$ is the outcome of an action the deployed system did
\emph{not} take, recovered by replaying a historical window rather than by modelling it. No
unconfoundedness assumption is invoked, because both $F(X_i^{a_0})$ and $F(X_i^{a})$ are
evaluated on the same fully observed historical window, so every recorded loss is a
realised measurement rather than an estimate of a latent potential outcome.

The central difficulty is that the quantity we care about depends on $Y$ and therefore
cannot be observed at deployment. Our remedy is to obtain it from the past. For a
historical window whose forecast target is already part of the legally visible history, the
same quantity is computable: the window is fully observed, its future is known, and it can
therefore be replayed as a \emph{pseudo-deployment}. Concretely, for a historical window $i$
we take its complete context $X_i$ and its already-realised future $y_i$, inject the
\emph{same} missingness mechanism that deployment uses (fixed before any method is run, see
\S\ref{sec:exp-setup}) to obtain $(X_i, M_i)$, and then actually execute every action
$a \in \calA$ with the frozen forecaster. The realised task loss of action $a$ on window $i$,
and the best achievable loss on that window, are
\begin{equation}
  L_{i,a} \;=\; \loss\!\left(F(T_a(X_i, M_i)),\, y_i\right),
  \qquad
  L_i^{\star} \;=\; \min_{b \in \calA} L_{i,b} ,
  \label{eq:loss}
\end{equation}
so that the \emph{action regret} of $a$ on window $i$ is
\begin{equation}
  r_{i,a} \;=\; L_{i,a} - L_i^{\star} \;\ge\; 0 .
  \label{eq:regret}
\end{equation}
Each executed pseudo-deployment contributes a record
\begin{equation}
  \calB \;=\; \bigl\{\, (z_i,\, a,\, L_{i,a}) \,\bigr\},
  \label{eq:replay-bank}
\end{equation}
where $z_i$ is the \emph{observable} state of the pseudo-deployment
(\S\ref{sec:method-scoring}). We also retain the relative utility
$g_{i,a} = L_{i,a_0} - L_{i,a}$, which reads naturally when a quantity is compared against
\textsc{Keep}, but the decision evidence is expressed through $L_{i,a}$ and $r_{i,a}$,
because those are the quantities the final decision is actually about.

The bank records what really happened when each legal intervention was applied to a real
historical window under the deployment missingness mechanism. It is not a model of the data,
and $L_{i,a}$ is not a reconstruction score: it is the realised change in forecasting loss.
Because $y_i$ belongs to TRAIN history, no test-time label is ever consulted. The bank is
built once, offline, and is then frozen; deployment reads it but never writes to it.
Figure~\ref{fig:concept}(b) summarises this construction. The bank is a discrete grid rather
than a model of missingness, and the resulting scope of validity is discussed in
Appendix~\ref{app:scope}.

\subsection{Regret-Aware Intervention Scoring}
\label{sec:method-scoring}
The decision we have to make is discrete, so what should be optimised is not the accuracy
with which a vector of action gains is regressed, but the ordering of actions on which the
decision depends. The two objectives are not the same: a scorer can improve its point-wise
gain error while leaving the argmax unchanged, or change the argmax while its average error
worsens. We therefore formulate scoring directly over action comparisons.

The state $z$ used by the scorer describes the \emph{request} and the \emph{candidate
interventions}, never the forecaster's response to a non-selected candidate, and it is kept
deliberately small:
\begin{itemize}\itemsep1pt
  \item \textbf{Episode state}: missing ratio, missing-run structure, the distance from the
  nearest missing position to the forecast origin, the pattern indicator, and summary
  statistics of the visible context.
  \item \textbf{Intervention state}: for each candidate version $X^{a}$ relative to the
  reference $X^{a_0}$, the action identity, the fraction of changed positions, and how much
  of the change falls near the forecast origin. These describe a difference between two
  \emph{inputs}, never a difference between forecasts.
\end{itemize}
The complete feature list, the normalisation statistics, and the exact learner are given in
Appendix~\ref{app:features}. Nothing in $z$ is computed from the future, and no
non-selected candidate's forecast is executed to build it.

Given $\calB$, the scorer produces a value $S_a(z)$ for every legal action and the deployed
decision is
\begin{equation}
  a^{\star} \;=\; \arg\max_{a \in \calA} S_a(z) .
  \label{eq:score}
\end{equation}
We consider two instantiations of $S_a$, both aligned with the ordering rather than with
point-wise gain error, and we freeze exactly one of them on the TRAIN gate before any main
experiment is run. In the \emph{pairwise} instantiation each replay record contributes
weighted preferences
\begin{equation}
  y_{i,a,b} \;=\; \mathbf{1}[L_{i,a} < L_{i,b}],
  \qquad
  \omega_{i,a,b} \;=\; |L_{i,a} - L_{i,b}| ,
  \label{eq:pairwise}
\end{equation}
a comparison model $q_\theta(a \succ b \mid z_i)$ is fitted to them, and the action score
aggregates the fitted comparisons,
\begin{equation}
  S_a^{\mathrm{pair}}(z) \;=\; \sum_{b \neq a} q_\theta(a \succ b \mid z) .
  \label{eq:score-pair}
\end{equation}
In the \emph{ensemble} instantiation the pairwise score is combined with expected-regret
estimates from two lightweight regressors,
\begin{equation}
  S_a(z) \;=\; \lambda_p S_a^{\mathrm{pair}}(z) \;-\; \lambda_r \hat r_a^{\mathrm{ridge}}(z)
  \;-\; \lambda_t \hat r_a^{\mathrm{tree}}(z) .
  \label{eq:score-ensemble}
\end{equation}
Both instantiations use the same replay bank, the same state $z$, and the same decision rule
\eqref{eq:score}; they differ only in how $\calB$ is converted into a ranking. Which one is
used, together with its fitted hyperparameters and the measured ordering agreement, is
reported in \S\ref{sec:exp-main} and \S\ref{sec:exp-ablation}.
% FREEZE_AFTER_R2_SELECTOR: 实验 Agent 冻结最终 selector 后，删除未采用的那一支（或降级为
% 消融项），并把采用支写实：$q_\theta$ 的模型族与容量、$\lambda_p,\lambda_r,\lambda_t$、
% 正则与早停口径、以及 $\tau$ 的选取位置。在此之前不得把任一支写成"我们的方法就是"。

\subsection{Selective Decision}
\label{sec:method-decision}
\textsc{Keep} is the reference action $a_0$ and competes in \eqref{eq:score} on the same
footing as every other legal action. Abstention is therefore not a separate guard bolted onto
a repair module: it is the outcome of the ranking whenever the untouched input is the
highest-scoring option. This distinction matters for what the method may claim. A
formulation that first selects a candidate and then applies a conservative threshold does
lower the harmful-intervention rate, but it does so by suppressing actions whose expected
regret is genuinely negative as well as actions whose evidence is merely thin; we document
that trade-off, measured in the preceding development cycle, in
Appendix~\ref{app:development}. The formulation used here instead lets the ranking decide, so
a wrongly suppressed intervention and a wrongly executed one are the same kind of error.

If the frozen scorer exposes a calibrated comparison against \textsc{Keep}, we add a
lightweight ambiguity check and fall back to the reference action,
\begin{equation}
  P(a^{\star} \succ a_0 \mid z) < \tau \;\Rightarrow\; a^{\star} = a_0 ,
  \label{eq:abstain}
\end{equation}
with $\tau$ selected on the TRAIN gate together with the other hyperparameters.
\eqref{eq:abstain} is ambiguity handling, not a guarantee: we do not claim coverage,
calibration, or any lower bound on the realised utility of a selected action, and every
decision-quality number we report is measured rather than bounded.

\subsection{Deployment and Complexity}
\label{sec:method-deploy}
A request is served in five steps: read the visible context and mask; compute the observable
state $z$; materialise the legal candidate inputs; score every candidate against the frozen
bank; execute $a^{\star}$ with the frozen forecaster and return the resulting forecast. When
$a^{\star} = a_0$ the input is untouched and the returned forecast is the reference one.
Under the frozen state definition above the scorer reads no forecast: $z$ is computed from
the input and the mask alone, so the frozen forecaster is executed once per request, for the
action the ranking selects (Appendix~\ref{app:calls}). There is no ensemble, no output
correction, no residual model, and no parameter update. Cost accounting is kept explicit and
separated: candidate materialisation, frozen-model calls, selector overhead, and total
request wall-clock are recorded individually, and no component is booked as free because a
cached artefact happened to exist.

% ===========================================================================
\section{Experiments}
\label{sec:exp}

\subsection{Experimental Setup}
\label{sec:exp-setup}
\paragraph{Datasets and backbones.}
The main matrix uses \ph{N_DATASETS} public forecasting sources: ETTh1, ETTh2, ETTm1,
ETTm2, Electricity~\citep{uci_electricity}, Exchange~\citep{lai2018lstnet},
Traffic~\citep{pems}, and Weather~\citep{wu2021autoformer}. Three frozen \tsfm{} families
are used: Chronos-Bolt~\citep{chronosbolt2025} (\emph{Bolt}),
TimesFM-2.5~\citep{das2024timesfm} (\emph{TimesFM}), and Chronos-2~\citep{chronos2_2025}
(\emph{Chronos-2}). Bolt and TimesFM are used for method development; Chronos-2 enters
only after the configuration is frozen, using the identical algorithm and hyperparameters.
This tests transfer across frozen \tsfm{} families rather than per-backbone tuning. Context
length is $L = 512$, main horizons are $H \in \{96, 192\}$, and longer horizons appear
only in Appendix~\ref{app:horizons}.

\paragraph{Governance actions.}
The catalog is fixed at five legal actions,
$\calA = \{\textsc{Keep},\ \textsc{Ffill},\ \textsc{Single TS-ICL},\ \textsc{Multi TS-ICL},\ \textsc{Context Ridge}\}$,
with \textsc{Keep} as the reference $a_0$. A complete and valid input is always mapped to
\textsc{Keep}; unsupported, aliased, and failed actions are recorded rather than silently
dropped.

\paragraph{Missingness protocol.}
Four deterministic patterns are frozen before any method runs: \textbf{P1 point missing};
\textbf{P2 target block}; \textbf{P3 shared block}; and \textbf{P4 tail missing}. Every mask
is derived by a deterministic hash of \texttt{source + parent + origin + horizon + pattern
+ protocol\_seed} and is never moved after forecasts are seen. The main experiment uses
$10\%$ severity; \S\ref{sec:exp-robust} adds $30\%$ and $50\%$.

\paragraph{Splits.}
Only legal TRAIN parents are used. Within TRAIN, parents are sorted by source and time and
divided into \emph{replay-fit} ($\sim$60\%, to build the replay bank and fit feature
scaling), \emph{gate} ($\sim$20\%, to select the selector hyperparameters and any legal
baseline hyperparameter), and \emph{TRAIN-eval} ($\sim$20\%, the final internal acceptance
check). Purging is enforced at $L + H$ and no parent's masks, patterns, or horizons cross a
split boundary. Replay pseudo-deployment severities are drawn from the same three levels by
deterministic hash, so the method is a single configuration that is not retuned per severity.
\S\ref{sec:exp-support} reuses this split rule while restricting how much of replay-fit is
available.

\paragraph{Baselines and controls.}
Five published baselines are fixed in advance and never swapped after results are seen:
TOI~\citep{wang2024taskoriented}, TOI-VSF~\citep{hao2025toivsf},
GIMCC~\citep{hao2025gimcc}, SRDI~\citep{xu2026srdi}, and
ChannelTokenFormer~\citep{jang2026channeltokenformer}. All five appear in every
non-ablation experiment in the main text. The four data-side methods receive a per-backbone
column; ChannelTokenFormer is a complete forecasting architecture trained for the task rather
than a data-side adapter, so its row in Table~\ref{tab:main} has no per-frozen-backbone
column and no number is fabricated for the missing cells. VIDA~\citep{liang2025vida} is held
back from the core set and reported as an extended comparison in
Appendix~\ref{app:extended}. Three controls accompany the five: \textsc{Native Keep},
\textsc{Best Fixed}, and \textbf{R2-CART}; a \emph{Catalog Oracle} using the future target is
reported in grey italics as a diagnostic upper reference and is excluded from all ranks.
Trainable baselines are fitted on TRAIN only.

\paragraph{Metrics, aggregation, and statistics.}
Four forecasting metrics are reported throughout: MASE~\citep{hyndman2006mase} as the primary
metric, plus MAE, RMSE, and MSE. MASE is averaged source-macro: first over parents, then over
sources, with the eight sources equally weighted. MAPE, CRPS, and QLIKE are not headline
metrics (MAPE is unstable near zero, CRPS is not exposed by every backbone or baseline, and
QLIKE is deferred to Appendix~\ref{app:finance}). Efficiency is reported separately and
never folded into an accuracy score. The primary statistical unit is the parent: paired
cluster bootstrap with \ph{N_BOOTSTRAP} resamples over parent blocks gives 95\% confidence
intervals on source-macro MASE differences against each baseline, with Holm
correction~\citep{holm1979simple} across baselines, and additionally we report average rank
and the number of won cells over source $\times$ backbone $\times$ horizon $\times$ pattern.
Mask variants of the same parent are never treated as independent samples.

\subsection{Main Results}
\label{sec:exp-main}
\emph{Can \introact{} improve forecasting under incomplete inputs across datasets and frozen
\tsfm{} backbones, compared with published methods designed for task-oriented or
incomplete-input forecasting?}

Table~\ref{tab:main} reports the main comparison at $10\%$ severity, over all eight sources,
three backbones, two horizons, and four patterns. \ph{MAIN_READING}.
\ph{MAIN_BACKBONE_CONSISTENCY}. \ph{MAIN_BASELINE_SPREAD}. \ph{MAIN_R2CART}. \ph{MAIN_RANK}.
\ph{MAIN_SELECTOR_CHOICE}. Per-source MAE and RMSE tables are deferred to
Appendix~\ref{app:datasets}. Deployment cost is reported separately in
Appendix~\ref{app:calls} and is never folded into an accuracy score: offline fitting is a
one-off cost, online cost covers candidate construction, frozen-model calls, and end-to-end
request latency, cached historical artefacts are not booked as zero, and cold start is
reported separately from hot latency.

\begin{table}[t]
  \caption{Main results at $10\%$ missingness. MASE is source-macro averaged and lower is
  better everywhere. The \emph{Type} column separates data-side methods, which can be
  composed with a frozen backbone, from ChannelTokenFormer, which is a complete forecasting
  architecture trained for the task: it has no per-frozen-backbone column, and its row is
  reported on the task-level comparison only, with no number fabricated for the missing
  cells. \best{Bold} is the best deployable method, \second{underline} the second best. The
  oracle row uses future targets for diagnosis only, is not deployable, and is excluded from
  ranks and from all claims.}
  \label{tab:main}
  \centering
  \small
  \setlength{\tabcolsep}{3.5pt}
\resizebox{\textwidth}{!}{\begin{tabular}{llccccccc}
    \toprule
    Method & Type & Bolt MASE $\downarrow$ & TimesFM MASE $\downarrow$
           & Chronos-2 MASE $\downarrow$ & Overall MASE $\downarrow$
           & MAE $\downarrow$ & RMSE $\downarrow$ & Rank $\downarrow$ \\
    \midrule
    Native KEEP        & Frozen TSFM            & \ph{KEEP_BOLT}   & \ph{KEEP_TF}   & \ph{KEEP_CH2}   & \ph{KEEP_OVR}   & \ph{KEEP_MAE}   & \ph{KEEP_RMSE}   & \ph{KEEP_RANK} \\
    Best Fixed         & Frozen TSFM            & \ph{BF_BOLT}     & \ph{BF_TF}     & \ph{BF_CH2}     & \ph{BF_OVR}     & \ph{BF_MAE}     & \ph{BF_RMSE}     & \ph{BF_RANK} \\
    R2-CART            & Selector               & \ph{R2_BOLT}     & \ph{R2_TF}     & \ph{R2_CH2}     & \ph{R2_OVR}     & \ph{R2_MAE}     & \ph{R2_RMSE}     & \ph{R2_RANK} \\
    TOI                & Task-oriented          & \ph{TOI_BOLT}    & \ph{TOI_TF}    & \ph{TOI_CH2}    & \ph{TOI_OVR}    & \ph{TOI_MAE}    & \ph{TOI_RMSE}    & \ph{TOI_RANK} \\
    TOI-VSF            & VSF                    & \ph{TOIVSF_BOLT} & \ph{TOIVSF_TF} & \ph{TOIVSF_CH2} & \ph{TOIVSF_OVR} & \ph{TOIVSF_MAE} & \ph{TOIVSF_RMSE} & \ph{TOIVSF_RANK} \\
    GIMCC              & VSF                    & \ph{GIMCC_BOLT}  & \ph{GIMCC_TF}  & \ph{GIMCC_CH2}  & \ph{GIMCC_OVR}  & \ph{GIMCC_MAE}  & \ph{GIMCC_RMSE}  & \ph{GIMCC_RANK} \\
    SRDI               & VSF                    & \ph{SRDI_BOLT}   & \ph{SRDI_TF}   & \ph{SRDI_CH2}   & \ph{SRDI_OVR}   & \ph{SRDI_MAE}   & \ph{SRDI_RMSE}   & \ph{SRDI_RANK} \\
    ChannelTokenFormer & Forecast model         & N/A\textsuperscript{*} & N/A\textsuperscript{*} & N/A\textsuperscript{*} & \ph{CTF_OVR} & \ph{CTF_MAE} & \ph{CTF_RMSE} & \ph{CTF_RANK} \\
    \midrule
    \rowcolor{bestgray}\introact{} & Frozen TSFM governance & \ph{OURS_BOLT} & \ph{OURS_TF} & \ph{OURS_CH2} & \ph{OURS_OVR} & \ph{OURS_MAE} & \ph{OURS_RMSE} & \ph{OURS_RANK} \\
    \midrule
    Catalog Oracle     & Diagnostic             & \oracle{\ph{ORACLE_BOLT}} & \oracle{\ph{ORACLE_TF}} & \oracle{\ph{ORACLE_CH2}} & \oracle{\ph{ORACLE_OVR}} & \oracle{\ph{ORACLE_MAE}} & \oracle{\ph{ORACLE_RMSE}} & \oracle{--} \\
    \bottomrule
  \end{tabular}}
  \vspace{-2pt}
  {\footnotesize\textsuperscript{*}Not applicable: ChannelTokenFormer is a trained forecasting
  architecture rather than a data-side adapter over a frozen backbone.}
\end{table}

\subsection{When Does Action Choice Matter?}
\label{sec:exp-opportunity}
\emph{Does a single fixed repair policy suffice, or does the useful action depend on the
episode?}

This experiment tests the premise of selective governance on the evaluation episodes. For
each parent we define the oracle opportunity of the episode,
\begin{equation}
  \Delta_i \;=\; L_{\textsc{Keep},i} - L_i^{\star},
  \qquad
  L_i^{\star} = \min_{a \in \calA} L_{i,a},
  \label{eq:opportunity}
\end{equation}
which is how much the best legal action could have improved on the untouched input. Because
$\Delta_i$ uses the future target, it is a diagnostic, not a deployable feature: the
\emph{strata} are defined by thresholds fitted on TRAIN only and then applied unchanged to
the evaluation episodes, giving \emph{no-op} (the reference input is already best, up to a
tolerance fixed on TRAIN), \emph{low opportunity}, and \emph{high opportunity}. No evaluation
episode is used to choose the thresholds.

All five published baselines enter, and Table~\ref{tab:opportunity} reports MASE within each
stratum under the identical partition. The table answers two questions at once: does
\introact{} avoid governance where nothing can be gained, and does it exploit action
heterogeneity where it exists? A method that is better only in the high-opportunity stratum
while degrading the no-op stratum has not solved the selection problem this paper poses.
\ph{OPPORTUNITY_READING}. \ph{OPPORTUNITY_NOOP}. \ph{OPPORTUNITY_HIGH}. The distribution of
$\Delta_i$ and the sensitivity of the stratum boundaries are reported in
Appendix~\ref{app:opportunity}.

\begin{table}[t]
  \caption{When does action choice matter? Episodes are partitioned into no-op,
  low-opportunity, and high-opportunity strata by TRAIN-only thresholds on the oracle
  opportunity $\Delta_i$ of \eqref{eq:opportunity}; MASE is source-macro averaged within each
  stratum and lower is better. The identical partition is applied to every method.}
  \label{tab:opportunity}
  \centering
  \small
  \setlength{\tabcolsep}{4pt}
  \begin{tabular}{lcccc}
    \toprule
    Method & No-op MASE $\downarrow$ & Low opportunity $\downarrow$ & High opportunity $\downarrow$ & Overall $\downarrow$ \\
    \midrule
    TOI                & \ph{NOOP_TOI}    & \ph{LOWOP_TOI}    & \ph{HIGHOP_TOI}    & \ph{OVR_TOI} \\
    TOI-VSF            & \ph{NOOP_TOIVSF} & \ph{LOWOP_TOIVSF} & \ph{HIGHOP_TOIVSF} & \ph{OVR_TOIVSF} \\
    GIMCC              & \ph{NOOP_GIMCC}  & \ph{LOWOP_GIMCC}  & \ph{HIGHOP_GIMCC}  & \ph{OVR_GIMCC} \\
    SRDI               & \ph{NOOP_SRDI}   & \ph{LOWOP_SRDI}   & \ph{HIGHOP_SRDI}   & \ph{OVR_SRDI} \\
    ChannelTokenFormer & \ph{NOOP_CTF}    & \ph{LOWOP_CTF}    & \ph{HIGHOP_CTF}    & \ph{OVR_CTF} \\
    \midrule
    \rowcolor{bestgray}\introact{} & \ph{NOOP_OURS} & \ph{LOWOP_OURS} & \ph{HIGHOP_OURS} & \ph{OVR_OURS} \\
    \bottomrule
  \end{tabular}
\end{table}

Figure~\ref{fig:opportunity} in Appendix~\ref{app:opportunity} visualises the stratum shares
together with the per-stratum improvement over \textsc{Keep}; the panel is a placeholder until
the main experiment produces the underlying CSV.

\subsection{Robustness to Missingness}
\label{sec:exp-robust}
\emph{Does the same frozen governance policy remain useful as missingness becomes more
severe and structurally different?}

The same five published baselines enter, and no method is retuned between severities.
Severities are $10\%$, $30\%$ and $50\%$, crossed with the four patterns and the three
backbones, with the selector configuration carried over unchanged from the TRAIN gate.
Table~\ref{tab:robust} reports the severity trend, the worst pattern cell, and the average
rank; the complete pattern $\times$ severity breakdown is in
Appendix~\ref{app:robustness}, and Figure~\ref{fig:robust} in
Appendix~\ref{app:robust-fig} shows the trend. \ph{ROBUST_READING}. \ph{ROBUST_50}.
\ph{ROBUST_WORST_PATTERN}.

\begin{table}[t]
  \caption{Robustness to missingness severity, with no method retuned between levels. The
  worst-pattern column is the largest degradation relative to \textsc{Native Keep} over all
  pattern and backbone cells at that severity; a negative entry means a method can be worse
  than doing nothing. Average rank is over method $\times$ severity $\times$ pattern $\times$
  backbone cells.}
  \label{tab:robust}
  \centering
  \small
  \setlength{\tabcolsep}{3pt}
\resizebox{\textwidth}{!}{\begin{tabular}{lcccccc}
    \toprule
    Method & 10\% MASE $\downarrow$ & 30\% MASE $\downarrow$ & 50\% MASE $\downarrow$ & Worst pattern $\downarrow$ & Avg.\ rank $\downarrow$ \\
    \midrule
    TOI                & \ph{TOI_10}    & \ph{TOI_30}    & \ph{TOI_50}    & \ph{TOI_WORST}    & \ph{TOI_ARANK} \\
    TOI-VSF            & \ph{TOIVSF_10} & \ph{TOIVSF_30} & \ph{TOIVSF_50} & \ph{TOIVSF_WORST} & \ph{TOIVSF_ARANK} \\
    GIMCC              & \ph{GIMCC_10}  & \ph{GIMCC_30}  & \ph{GIMCC_50}  & \ph{GIMCC_WORST}  & \ph{GIMCC_ARANK} \\
    SRDI               & \ph{SRDI_10}   & \ph{SRDI_30}   & \ph{SRDI_50}   & \ph{SRDI_WORST}   & \ph{SRDI_ARANK} \\
    ChannelTokenFormer & \ph{CTF_10}    & \ph{CTF_30}    & \ph{CTF_50}    & \ph{CTF_WORST}    & \ph{CTF_ARANK} \\
    \midrule
    \rowcolor{bestgray}\introact{} & \ph{OURS_10} & \ph{OURS_30} & \ph{OURS_50} & \ph{OURS_WORST} & \ph{OURS_ARANK} \\
    \bottomrule
  \end{tabular}}
\end{table}

\subsection{Historical Support Efficiency}
\label{sec:exp-support}
\emph{Does counterfactual replay provide an effective way to exploit limited historical
supervision for incomplete-input decision making?}

Replay supervision is a bounded resource: it is limited by how many historical windows are
legally available, not by compute. This experiment restricts the historical support that every
method may use to $25\%$, $50\%$, and $100\%$ of the replay-fit portion of TRAIN under the
identical split rule of \S\ref{sec:exp-setup}, and re-fits each method from scratch at each
level. All five published baselines enter. The compact reading statement is below and the
full source-macro table with the offline-cost column is deferred to
Appendix~\ref{app:support}: \ph{SUPPORT_READING}. \ph{SUPPORT_25}, \ph{SUPPORT_50},
\ph{SUPPORT_100}, \ph{SUPPORT_OFFLINE_COST}.

\subsection{Ablation Study}
\label{sec:exp-ablation}
\emph{Which component carries the gain, and does regret-aware ranking actually matter?}

The ablation is fixed at six rows and reuses the unified main protocol: A1 replaces
regret-aware ranking with direct gain regression; A2 drops the pairwise weights
$\omega_{i,a,b}$ of \eqref{eq:pairwise}; A3 drops the intervention state, A4 drops the
episode context; A5 removes \textsc{Keep} from the decision space so the system must
intervene on every request. The full six-row table is deferred to
Appendix~\ref{app:ablation-details}; the headline reading is \ph{ABLATION_READING}.

% ===========================================================================
\section{Conclusion}
\label{sec:conclusion}
We formulated incomplete-input handling for frozen time-series foundation models as
selective data governance. It is an \emph{intervention decision problem}, where a feasible
repair is not automatically beneficial and the untouched input is a legitimate answer;
historical replay provides \emph{task-level supervision}, since the realised consequence of an
intervention is measured rather than modelled; and \emph{decision-aligned selection}
determines when to intervene without modifying the frozen backbone, returning a forecast from
exactly one executed input version. \ph{CONCLUSION_CLOSING}.

\subsubsection*{AI use statement}

\subsubsection*{AI use statement}
In this work, we used generative AI tools for drafting and editing prose, for
implementing and refactoring experiment code, and for producing figure sources from
explicit geometric specifications. We have not used generative AI tools to generate
experimental results, to select or curate datasets, or to make methodological decisions.
All AI-assisted code was executed and verified against frozen artefacts, and every
numerical claim in this paper is produced by the experiment pipeline described in
Appendix~\ref{app:reprodetails} and is traceable to a recorded artefact. We take
responsibility for the final content of this work, including text, claims, and artefacts
produced with the aid of generative AI.
\subsubsection*{Ethics statement}
This work studies data governance for frozen forecasting models on public benchmark
datasets. It does not involve human subjects, personal data, or model deployment
decisions with direct societal consequences. The governance actions considered here never
overwrite valid observations, which is a safeguard by construction: an intervention that
silently replaced real measurements would be an integrity problem regardless of its
forecasting effect. We report negative results and failed configurations rather than
removing them, and we do not claim safety, robustness, or generality beyond the measured
conditions.

\subsubsection*{Reproducibility statement}
All datasets used are public benchmarks. The governance action catalog, the missingness
protocol, the mask derivation rule, the TRAIN-only split with $L+H$ purging, the state
feature list, the hyperparameter grid, and the failure policy are fixed in configuration
files that are frozen before the main experiment and identified by commit hash. Every
result row records the source, parent, origin, horizon, pattern, severity, mask hash,
backbone revision, method, action, input hash, prediction hash, metric values, runtime
components, failure state, code SHA, and resolved configuration, so that any reported
number can be traced back to the artefact that produced it. Aggregate-only reporting is
avoided: raw predictions are retained. Details are in Appendix~\ref{app:reprodetails}.

\bibliography{references}
\bibliographystyle{iclr2027_conference}

% ===========================================================================
% ===========================================================================
% 附录（由正文下沉 + 完整化）。所有未运行数字一律 \ph{KEY}，禁止编造。
% ===========================================================================



\appendix

\newpage



\section*{Appendix map}
\label{app:map}
Table~\ref{tab:app-map} lists what each appendix section contains, so that a reader looking
for a specific piece of evidence does not have to scan the whole appendix. Sections~A--D fix
the protocol and the method details, E--I report the evidence behind the main-text tables,
J records what was tried and rejected, and K--M cover the boundary regimes.

\begin{table}[h]
  \caption{Appendix map. Sections are lettered in order of appearance; this table is
  unnumbered and does not consume a section letter.}
  \label{tab:app-map}
  \centering
  \small
  \setlength{\tabcolsep}{4pt}
\resizebox{\textwidth}{!}{\begin{tabular}{lll}
    \toprule
    Section & Contents & Question it answers \\
    \midrule
    A. Scope and assumptions & replay-bank scope, closed catalog, determinism & what is out of scope, and which claims are therefore not made? \\
    B. Full method details & action catalog, observable state, relation to OPE and selective prediction & what exactly is the method, and how is it positioned? \\
    C. Replay construction & splits and purging, bank construction, bank size & how was the replay bank built, and how much of it is needed? \\
    D. Dataset and missingness protocol & the four frozen patterns and the mask derivation rule & what missingness is being simulated, and how is it made deterministic? \\
    E. Full main results & per-source tables, reconstruction-vs-utility, governance diagnostics, opportunity strata, extended comparison & does the aggregate hold source by source and stratum by stratum? \\
    F. Robustness breakdown & $10/30/50\%$ tables, per-pattern results, severity trend & does the frozen policy degrade slower than the baselines? \\
    G. Historical support results & every method re-fitted at $25/50/100\%$ support & how much history does each method need? \\
    H. Ablation details & full outcome vector and ordering stability & which component carries the gain? \\
    I. Runtime and failure audit & latency, call counts, failures, timeouts, cold start & what does a request actually cost? \\
    J. Development negative results & local gain averaging with lower-confidence gating, gain regression, the sealed v4.4 TRAIN-eval batch & what was tried and rejected, and why? \\
    K. Complete-input contract & required behaviour by input regime & what happens when nothing is missing? \\
    L. Financial case study & finance parents, QLIKE & what happens outside the benchmark sources? \\
    M. Additional seeds and sensitivity & mask-seed stability, extra horizons, record schema & is the conclusion an artefact of one mask, or of the chosen horizons? \\
    \bottomrule
  \end{tabular}}
\end{table}


\section{Scope and Assumptions}
\label{app:scope}
This section states what the replay bank is, what it is not, and which claims the paper
therefore does and does not make.

The bank is a discrete grid (four patterns $\times$ three severities $\times$ registered
sources), not a continuous model of missingness. If a request falls outside that grid --- a
different missingness mechanism, a different missing-rate regime, or an unrepresented source
family --- the scorer still returns a ranking, but the retrieved historical losses need not
transfer. We therefore read \S\ref{sec:exp-robust} as a within-grid check rather than as
out-of-grid generalisation, and the finance case study (Appendix~\ref{app:finance}) is the
only contact with a non-synthetic gap in this paper.

Three further assumptions are worth stating explicitly. First, the action catalog is closed and
given: we do not search for new interventions, and a wider catalog would change every number
in this paper without changing the decision problem. Second, the deployment missingness
mechanism is assumed to be the one that the replay injects; a request produced by a different
mechanism is handled by the same scorer but is outside the scope of the reported results.
Third, the frozen forecaster is assumed deterministic and revision-pinned: every replay record
is tied to the exact model revision that produced it, and a model update invalidates the bank
rather than silently degrading it.


\section{Full Method Details}
\label{app:method-details}
This section expands \S\ref{sec:method}: the closed action catalog, the observable state used
by the scorer, and the relation of the decision layer to off-policy evaluation and selective
prediction. The learner that turns replay records into a ranking is frozen on the TRAIN gate
and reported in \S\ref{sec:exp-main}; its implementation is not restated here because it is
selected by the protocol rather than fixed by the paper's formulation.

\subsection{Governance Action Catalog}
\label{app:catalog}
Table~\ref{tab:app-catalog} defines the five legal actions. The catalog is closed: no
action outside this set is ever executed, and no action may overwrite an observation that
has actually arrived.

\begin{table}[h]
  \caption{The five legal governance actions. ``Observed entries'' refers to entries that
  are present and valid in $X_i^{a_0}$; an action that modified them would be a data
  integrity violation regardless of its forecasting effect.}
  \label{tab:app-catalog}
  \centering
  \small
  \setlength{\tabcolsep}{4pt}
\resizebox{\textwidth}{!}{\begin{tabular}{llcc}
    \toprule
    Action & Input version produced & May modify observed entries? & Extra frozen-model calls \\
    \midrule
    \textsc{Keep} $a_0$   & the incomplete context, unchanged & --- (reference) & 0 \\
    \textsc{Ffill}        & forward-fill from the last valid value & no & 1 \\
    \textsc{Single TS-ICL}& univariate in-context repair of the target channel & no & 1 \\
    \textsc{Multi TS-ICL} & multivariate in-context repair using visible channels & no & 1 \\
    \textsc{Context Ridge}& ridge reconstruction from the visible context & no & 1 \\
    \bottomrule
  \end{tabular}}
\end{table}

A complete and valid input is always mapped to \textsc{Keep} and the reference forecast is
returned unchanged. Unsupported, aliased, and failed actions are recorded with an explicit
status; they are never silently dropped and never counted as successes.

\subsection{Observable State Features}
\label{app:features}
The state $z_i$ is lightweight and fully observable at request time. It contains no future
target information and, under the frozen state definition of \S\ref{sec:method-scoring}, no
quantity computed from the frozen model's output at all: scoring happens before any forecast
is produced. Table~\ref{tab:app-features} lists the two groups.

\begin{table}[h]
  \caption{The two observable state groups. Every feature is computable from the request
  alone; none requires the future target, and none requires a forecast. Group~I is evaluated
  once per candidate action; group~E is evaluated once per request.}
  \label{tab:app-features}
  \centering
  \small
  \setlength{\tabcolsep}{4pt}
\begin{tabular}{@{}l p{0.60\textwidth} c@{}}
    \toprule
    Group & Features & Dim. \\
    \midrule
    E Episode state      & missing ratio, number of missing runs, longest run length, distance of the missing block to the origin, pattern indicator, and summary statistics of the visible context (level, local trend, seasonality strength, volatility) & \ph{DIM_E} \\
    I Intervention state & action identity, fraction of entries changed relative to the reference input, and the share of the change that falls near the forecast origin & \ph{DIM_I} \\
    \bottomrule
  \end{tabular}
\end{table}

Features are standardised using statistics fitted on replay-fit parents only. The normalisation
is part of the frozen configuration and is not refitted on the gate or the evaluation splits.
The exact learner that maps $z$ to a ranking, together with its regularisation and early
stopping, is frozen on the gate split and reported in \S\ref{sec:exp-main}; it is not restated
here because it is selected by the protocol rather than fixed by the paper's formulation.


\subsection{Relation to Off-Policy Evaluation and Selective Prediction}
\label{app:ope-positioning}
Our decision layer is formally close to conservative policy selection in contextual
bandits~\citep{abbasiyadkori2011improved,lattimore2020bandit} and to off-policy
evaluation~\citep{strehl2010learning,dudik2011doubly,swaminathan2015counterfactual,
jiang2016doubly,thomas2016data}, but three differences matter here. First, an arm's value is
not a designed reward but the realised loss of a frozen, non-differentiable forecaster, so
there is no reward model to fit and no simulator to query: the only legal supervision is a
replayed historical window, tied to the exact model revision that produced it. Second, the
untouched input is one of the arms rather than a fallback, so a wrong decision is asymmetric
and directly measurable, which is why we report harmful loss and action-opportunity strata
rather than regret alone. Third, the exploration budget is zero: no action is ever executed at
deployment in order to learn, so the bank is a fixed historical artefact and the method has no
online update.

Selective prediction and learning to defer~\citep{chow1970optimum,vovk2005algorithmic,
geifman2017selective,madras2018predict,mozannar2020consistent} are the closest relatives, but
they decide whether to \emph{predict} or to hand the case to a human or a second model. We
decide whether to \emph{modify the input} of one fixed predictor, and our abstention returns
that predictor's own reference forecast. Work on sequential evidence acquisition and decision
value~\citep{he2024dime,kang2020forecastwithforecasts} formalises when information is worth
its cost; we make no information-theoretic optimality claim and report the full cost
accounting separately (Appendix~\ref{app:calls}).

\section{Replay Construction}
\label{app:replay}
This section records how the replay bank is built and how much of it is needed. Nothing here
is a result; it is the construction that the result tables are read against.

\subsection{Splits, Purging, and Replay-Bank Construction}
\label{app:splits}
Sources are split at the parent level only. Within each source the registered legal TRAIN
parents are sorted by origin time and partitioned into \emph{replay-fit}, \emph{gate}, and
\emph{TRAIN-eval} at approximately $60/20/20$. The replay bank is built exclusively from
replay-fit parents; the selector hyperparameters and any legal baseline hyperparameter are selected on the
gate split; TRAIN-eval is the final internal acceptance check before the configuration is
frozen. Purging is enforced at $L + H$, so no parent's context window overlaps a later
split's target window. Table~\ref{tab:app-splits} records the realised sizes.

\begin{table}[h]
  \caption{Realised TRAIN-only split sizes and parent capacities. Counts are of legal
  parents, not of correlated mask variants; a parent with several patterns or severities
  still contributes one unit to every statistic in the paper.}
  \label{tab:app-splits}
  \centering
  \small
  \setlength{\tabcolsep}{4pt}
\resizebox{\textwidth}{!}{\begin{tabular}{lcccc}
    \toprule
    Split & Parents & Origins & Sources covered & Used for \\
    \midrule
    Replay-fit   & \ph{FIT_PARENTS}   & \ph{FIT_ORIGINS}   & \ph{FIT_SOURCES}   & replay bank, baseline training, feature normalisation \\
    Gate         & \ph{GATE_PARENTS}  & \ph{GATE_ORIGINS}  & \ph{GATE_SOURCES}  & selector and baseline hyperparameters \\
    TRAIN-eval   & \ph{EVAL_PARENTS}  & \ph{EVAL_ORIGINS}  & \ph{EVAL_SOURCES}  & frozen-configuration acceptance check \\
    \midrule
    Total legal TRAIN & \ph{TRAIN_PARENTS} & \ph{TRAIN_ORIGINS} & \ph{N_DATASETS} & --- \\
    \bottomrule
  \end{tabular}}
\end{table}

The replay bank is constructed as follows. For every replay-fit parent we take the
historical window $(X_i, y_i)$ whose future is already legal history, inject the deployment
missingness mechanism to obtain the incomplete version $X_i^{a_0}$, materialise each legal
action $a \in \calA$ into $X_i^{a}$, actually execute the frozen forecaster on each
materialised version, and store the realised gain
\begin{equation}
  g_{i,a} \;=\; \loss\big(F(X_i^{a_0}),\, y_i\big) - \loss\big(F(X_i^{a}),\, y_i\big),
  \label{eq:app-gain}
\end{equation}
so that $g_{i,a} > 0$ means the action helped and $g_{i,a} < 0$ means it hurt. The bank is
$\calB = \{(z_i, a, g_{i,a})\}$, where $z_i$ is the observable task state defined in
Appendix~\ref{app:features}. Pseudo-deployment severities are drawn from
$\{10\%, 30\%, 50\%\}$ by deterministic hash of the same key that defines the mask, so the
bank covers the severity range without the method being retuned per severity. Because the
bank is built once from TRAIN and never rebuilt at deployment time, the offline cost is
paid once; Appendix~\ref{app:calls} separates it from the per-request online cost.

\subsection{Replay Bank Size}
\label{app:replaysize}
The replay bank is subsampled to $25\%$, $50\%$, and $100\%$ of its replay-fit support to
measure how much the method depends on the amount of historical evidence. This analysis is
not part of the main text and was not used to select any hyperparameter; the frozen
configuration is the one trained on the full bank. If performance at $25\%$ matches
performance at $100\%$, the honest reading is that the bank is already saturated at this
scale, not that the method is data-efficient in general.

\begin{table}[h]
  \caption{Replay-bank size ablation. The bank is subsampled by replay-fit parent, keeping
  every action and severity per retained parent, so the subsample changes the amount of
  history rather than its composition.}
  \label{tab:app-replaysize}
  \centering
  \small
  \setlength{\tabcolsep}{4pt}
  \begin{tabular}{lccccc}
    \toprule
    Bank fraction & MASE $\downarrow$ & Intervention Rate & HIR $\downarrow$ & Harmful Loss $\downarrow$ & Calls $\downarrow$ \\
    \midrule
    $25\%$  & \ph{RS25_MASE}  & \ph{RS25_IR}  & \ph{RS25_HIR}  & \ph{RS25_HL}  & \ph{RS25_CALLS} \\
    $50\%$  & \ph{RS50_MASE}  & \ph{RS50_IR}  & \ph{RS50_HIR}  & \ph{RS50_HL}  & \ph{RS50_CALLS} \\
    $100\%$ & \ph{RS100_MASE} & \ph{RS100_IR} & \ph{RS100_HIR} & \ph{RS100_HL} & \ph{RS100_CALLS} \\
    \bottomrule
  \end{tabular}
\end{table}

\section{Dataset and Missingness Protocol}
\label{app:data-protocol}
The missingness mechanism is frozen before any method is run, and every mask is derived
deterministically from a key that identifies the parent and the pattern. This section fixes
the four patterns and the derivation rule.

\subsection{Missingness Patterns}
\label{app:patterns}
Table~\ref{tab:app-patterns} defines the four deterministic patterns. Every mask is derived
by a deterministic hash of
\texttt{source + parent + origin + horizon + pattern + protocol\_seed} and is never moved
after forecasts are seen.

\begin{table}[h]
  \caption{The four frozen missingness patterns. All four are applied to the same legal
  parents, so pattern effects are measured on a common parent set rather than on different
  subsets of the data.}
  \label{tab:app-patterns}
  \centering
  \small
  \setlength{\tabcolsep}{4pt}
\begin{tabular}{@{}l p{0.34\textwidth} p{0.40\textwidth}@{}}
    \toprule
    Pattern & What is removed & Why it is included \\
    \midrule
    P1 point missing    & scattered single positions & the classical random-missing regime \\
    P2 target block     & a contiguous block on the target variable & a sensor outage on the variable being forecast \\
    P3 shared block     & several visible variables over the same interval & a channel or acquisition outage affecting the whole panel \\
    P4 tail missing     & a block adjacent to the forecast origin & observations that have not arrived yet at request time \\
    \bottomrule
  \end{tabular}
\end{table}

\section{Full Main Results}
\label{app:main-full}
This section expands \S\ref{sec:exp-main} and \S\ref{sec:exp-opportunity}: per-source tables,
the paired reconstruction-versus-utility table, the governance diagnostics figure, the
opportunity strata, and the extended comparison.

\subsection{Per-Source Results}
\label{app:datasets}
Tables~\ref{tab:app-src-bolt}--\ref{tab:app-src-ch2} give the complete per-source results at
$10\%$ severity for the three frozen backbones, all four metrics and both main horizons.
Rows are generated directly from the result group \texttt{v44\_main\_results} rather than
transcribed by hand, so the appendix and the main table cannot disagree.

\begin{table}[h]
  \caption{Full per-source results on the Bolt backbone at $10\%$ severity, all four metrics and both main horizons. Lower is better everywhere in this table.}
  \label{tab:app-src-bolt}
  \centering
  \small
  \setlength{\tabcolsep}{4pt}
\resizebox{\textwidth}{!}{\begin{tabular}{llcccc}
    \toprule
    Source & $H$ & MASE $\downarrow$ & MAE $\downarrow$ & RMSE $\downarrow$ & MSE $\downarrow$ \\
    \midrule
    ETTh1 & 96 & \ph{ETTH1-96-MASE} & \ph{ETTH1-96-MSE} & \ph{ETTH1-96-MAE} & \ph{ETTH1-96-RMSE} \\
     & 192 & \ph{ETTH1-192-MASE} & \ph{ETTH1-192-MSE} & \ph{ETTH1-192-MAE} & \ph{ETTH1-192-RMSE} \\
    \midrule
    ETTh2 & 96 & \ph{ETTH2-96-MASE} & \ph{ETTH2-96-MSE} & \ph{ETTH2-96-MAE} & \ph{ETTH2-96-RMSE} \\
     & 192 & \ph{ETTH2-192-MASE} & \ph{ETTH2-192-MSE} & \ph{ETTH2-192-MAE} & \ph{ETTH2-192-RMSE} \\
    \midrule
    ETTm1 & 96 & \ph{ETTM1-96-MASE} & \ph{ETTM1-96-MSE} & \ph{ETTM1-96-MAE} & \ph{ETTM1-96-RMSE} \\
     & 192 & \ph{ETTM1-192-MASE} & \ph{ETTM1-192-MSE} & \ph{ETTM1-192-MAE} & \ph{ETTM1-192-RMSE} \\
    \midrule
    ETTm2 & 96 & \ph{ETTM2-96-MASE} & \ph{ETTM2-96-MSE} & \ph{ETTM2-96-MAE} & \ph{ETTM2-96-RMSE} \\
     & 192 & \ph{ETTM2-192-MASE} & \ph{ETTM2-192-MSE} & \ph{ETTM2-192-MAE} & \ph{ETTM2-192-RMSE} \\
    \midrule
    Electricity & 96 & \ph{ELEC-96-MASE} & \ph{ELEC-96-MSE} & \ph{ELEC-96-MAE} & \ph{ELEC-96-RMSE} \\
     & 192 & \ph{ELEC-192-MASE} & \ph{ELEC-192-MSE} & \ph{ELEC-192-MAE} & \ph{ELEC-192-RMSE} \\
    \midrule
    Exchange & 96 & \ph{EXCH-96-MASE} & \ph{EXCH-96-MSE} & \ph{EXCH-96-MAE} & \ph{EXCH-96-RMSE} \\
     & 192 & \ph{EXCH-192-MASE} & \ph{EXCH-192-MSE} & \ph{EXCH-192-MAE} & \ph{EXCH-192-RMSE} \\
    \midrule
    Traffic & 96 & \ph{TRAF-96-MASE} & \ph{TRAF-96-MSE} & \ph{TRAF-96-MAE} & \ph{TRAF-96-RMSE} \\
     & 192 & \ph{TRAF-192-MASE} & \ph{TRAF-192-MSE} & \ph{TRAF-192-MAE} & \ph{TRAF-192-RMSE} \\
    \midrule
    Weather & 96 & \ph{WTHR-96-MASE} & \ph{WTHR-96-MSE} & \ph{WTHR-96-MAE} & \ph{WTHR-96-RMSE} \\
     & 192 & \ph{WTHR-192-MASE} & \ph{WTHR-192-MSE} & \ph{WTHR-192-MAE} & \ph{WTHR-192-RMSE} \\
    \midrule
    \rowcolor{bestgray}Source-macro & 96  & \ph{BOLT-96-MASE} & \ph{BOLT-96-MSE} & \ph{BOLT-96-MAE} & \ph{BOLT-96-RMSE} \\
    \rowcolor{bestgray}Source-macro & 192 & \ph{BOLT-192-MASE} & \ph{BOLT-192-MSE} & \ph{BOLT-192-MAE} & \ph{BOLT-192-RMSE} \\
    \bottomrule
  \end{tabular}}
\end{table}

The Bolt backbone is the reference for the main table, so its per-source rows are the ones the main table aggregates.

\begin{table}[h]
  \caption{Full per-source results on the TimesFM backbone at $10\%$ severity, all four metrics and both main horizons.}
  \label{tab:app-src-tf}
  \centering
  \small
  \setlength{\tabcolsep}{4pt}
\resizebox{\textwidth}{!}{\begin{tabular}{llcccc}
    \toprule
    Source & $H$ & MASE $\downarrow$ & MAE $\downarrow$ & RMSE $\downarrow$ & MSE $\downarrow$ \\
    \midrule
    ETTh1 & 96 & \ph{TF-ETTH1-96-MASE} & \ph{TF-ETTH1-96-MSE} & \ph{TF-ETTH1-96-MAE} & \ph{TF-ETTH1-96-RMSE} \\
     & 192 & \ph{TF-ETTH1-192-MASE} & \ph{TF-ETTH1-192-MSE} & \ph{TF-ETTH1-192-MAE} & \ph{TF-ETTH1-192-RMSE} \\
    \midrule
    ETTh2 & 96 & \ph{TF-ETTH2-96-MASE} & \ph{TF-ETTH2-96-MSE} & \ph{TF-ETTH2-96-MAE} & \ph{TF-ETTH2-96-RMSE} \\
     & 192 & \ph{TF-ETTH2-192-MASE} & \ph{TF-ETTH2-192-MSE} & \ph{TF-ETTH2-192-MAE} & \ph{TF-ETTH2-192-RMSE} \\
    \midrule
    ETTm1 & 96 & \ph{TF-ETTM1-96-MASE} & \ph{TF-ETTM1-96-MSE} & \ph{TF-ETTM1-96-MAE} & \ph{TF-ETTM1-96-RMSE} \\
     & 192 & \ph{TF-ETTM1-192-MASE} & \ph{TF-ETTM1-192-MSE} & \ph{TF-ETTM1-192-MAE} & \ph{TF-ETTM1-192-RMSE} \\
    \midrule
    ETTm2 & 96 & \ph{TF-ETTM2-96-MASE} & \ph{TF-ETTM2-96-MSE} & \ph{TF-ETTM2-96-MAE} & \ph{TF-ETTM2-96-RMSE} \\
     & 192 & \ph{TF-ETTM2-192-MASE} & \ph{TF-ETTM2-192-MSE} & \ph{TF-ETTM2-192-MAE} & \ph{TF-ETTM2-192-RMSE} \\
    \midrule
    Electricity & 96 & \ph{TF-ELEC-96-MASE} & \ph{TF-ELEC-96-MSE} & \ph{TF-ELEC-96-MAE} & \ph{TF-ELEC-96-RMSE} \\
     & 192 & \ph{TF-ELEC-192-MASE} & \ph{TF-ELEC-192-MSE} & \ph{TF-ELEC-192-MAE} & \ph{TF-ELEC-192-RMSE} \\
    \midrule
    Exchange & 96 & \ph{TF-EXCH-96-MASE} & \ph{TF-EXCH-96-MSE} & \ph{TF-EXCH-96-MAE} & \ph{TF-EXCH-96-RMSE} \\
     & 192 & \ph{TF-EXCH-192-MASE} & \ph{TF-EXCH-192-MSE} & \ph{TF-EXCH-192-MAE} & \ph{TF-EXCH-192-RMSE} \\
    \midrule
    Traffic & 96 & \ph{TF-TRAF-96-MASE} & \ph{TF-TRAF-96-MSE} & \ph{TF-TRAF-96-MAE} & \ph{TF-TRAF-96-RMSE} \\
     & 192 & \ph{TF-TRAF-192-MASE} & \ph{TF-TRAF-192-MSE} & \ph{TF-TRAF-192-MAE} & \ph{TF-TRAF-192-RMSE} \\
    \midrule
    Weather & 96 & \ph{TF-WTHR-96-MASE} & \ph{TF-WTHR-96-MSE} & \ph{TF-WTHR-96-MAE} & \ph{TF-WTHR-96-RMSE} \\
     & 192 & \ph{TF-WTHR-192-MASE} & \ph{TF-WTHR-192-MSE} & \ph{TF-WTHR-192-MAE} & \ph{TF-WTHR-192-RMSE} \\
    \midrule
    \rowcolor{bestgray}Source-macro & 96  & \ph{TF-96-MASE} & \ph{TF-96-MSE} & \ph{TF-96-MAE} & \ph{TF-96-RMSE} \\
    \rowcolor{bestgray}Source-macro & 192 & \ph{TF-192-MASE} & \ph{TF-192-MSE} & \ph{TF-192-MAE} & \ph{TF-192-RMSE} \\
    \bottomrule
  \end{tabular}}
\end{table}

TimesFM is used for method development alongside Bolt, so a per-source view here is the closest thing to a development-time breakdown.

\begin{table}[h]
  \caption{Full per-source results on the Chronos-2 backbone at $10\%$ severity, all four metrics and both main horizons. Chronos-2 takes no part in method design: the identical frozen configuration is applied to it after the fact.}
  \label{tab:app-src-ch2}
  \centering
  \small
  \setlength{\tabcolsep}{4pt}
\resizebox{\textwidth}{!}{\begin{tabular}{llcccc}
    \toprule
    Source & $H$ & MASE $\downarrow$ & MAE $\downarrow$ & RMSE $\downarrow$ & MSE $\downarrow$ \\
    \midrule
    ETTh1 & 96 & \ph{CH2-ETTH1-96-MASE} & \ph{CH2-ETTH1-96-MSE} & \ph{CH2-ETTH1-96-MAE} & \ph{CH2-ETTH1-96-RMSE} \\
     & 192 & \ph{CH2-ETTH1-192-MASE} & \ph{CH2-ETTH1-192-MSE} & \ph{CH2-ETTH1-192-MAE} & \ph{CH2-ETTH1-192-RMSE} \\
    \midrule
    ETTh2 & 96 & \ph{CH2-ETTH2-96-MASE} & \ph{CH2-ETTH2-96-MSE} & \ph{CH2-ETTH2-96-MAE} & \ph{CH2-ETTH2-96-RMSE} \\
     & 192 & \ph{CH2-ETTH2-192-MASE} & \ph{CH2-ETTH2-192-MSE} & \ph{CH2-ETTH2-192-MAE} & \ph{CH2-ETTH2-192-RMSE} \\
    \midrule
    ETTm1 & 96 & \ph{CH2-ETTM1-96-MASE} & \ph{CH2-ETTM1-96-MSE} & \ph{CH2-ETTM1-96-MAE} & \ph{CH2-ETTM1-96-RMSE} \\
     & 192 & \ph{CH2-ETTM1-192-MASE} & \ph{CH2-ETTM1-192-MSE} & \ph{CH2-ETTM1-192-MAE} & \ph{CH2-ETTM1-192-RMSE} \\
    \midrule
    ETTm2 & 96 & \ph{CH2-ETTM2-96-MASE} & \ph{CH2-ETTM2-96-MSE} & \ph{CH2-ETTM2-96-MAE} & \ph{CH2-ETTM2-96-RMSE} \\
     & 192 & \ph{CH2-ETTM2-192-MASE} & \ph{CH2-ETTM2-192-MSE} & \ph{CH2-ETTM2-192-MAE} & \ph{CH2-ETTM2-192-RMSE} \\
    \midrule
    Electricity & 96 & \ph{CH2-ELEC-96-MASE} & \ph{CH2-ELEC-96-MSE} & \ph{CH2-ELEC-96-MAE} & \ph{CH2-ELEC-96-RMSE} \\
     & 192 & \ph{CH2-ELEC-192-MASE} & \ph{CH2-ELEC-192-MSE} & \ph{CH2-ELEC-192-MAE} & \ph{CH2-ELEC-192-RMSE} \\
    \midrule
    Exchange & 96 & \ph{CH2-EXCH-96-MASE} & \ph{CH2-EXCH-96-MSE} & \ph{CH2-EXCH-96-MAE} & \ph{CH2-EXCH-96-RMSE} \\
     & 192 & \ph{CH2-EXCH-192-MASE} & \ph{CH2-EXCH-192-MSE} & \ph{CH2-EXCH-192-MAE} & \ph{CH2-EXCH-192-RMSE} \\
    \midrule
    Traffic & 96 & \ph{CH2-TRAF-96-MASE} & \ph{CH2-TRAF-96-MSE} & \ph{CH2-TRAF-96-MAE} & \ph{CH2-TRAF-96-RMSE} \\
     & 192 & \ph{CH2-TRAF-192-MASE} & \ph{CH2-TRAF-192-MSE} & \ph{CH2-TRAF-192-MAE} & \ph{CH2-TRAF-192-RMSE} \\
    \midrule
    Weather & 96 & \ph{CH2-WTHR-96-MASE} & \ph{CH2-WTHR-96-MSE} & \ph{CH2-WTHR-96-MAE} & \ph{CH2-WTHR-96-RMSE} \\
     & 192 & \ph{CH2-WTHR-192-MASE} & \ph{CH2-WTHR-192-MSE} & \ph{CH2-WTHR-192-MAE} & \ph{CH2-WTHR-192-RMSE} \\
    \midrule
    \rowcolor{bestgray}Source-macro & 96  & \ph{CH2-96-MASE} & \ph{CH2-96-MSE} & \ph{CH2-96-MAE} & \ph{CH2-96-RMSE} \\
    \rowcolor{bestgray}Source-macro & 192 & \ph{CH2-192-MASE} & \ph{CH2-192-MSE} & \ph{CH2-192-MAE} & \ph{CH2-192-RMSE} \\
    \bottomrule
  \end{tabular}}
\end{table}

Chronos-2 never participates in method design or hyperparameter selection; it is included to test whether the frozen configuration transfers to a family it was not developed on.

\subsection{Reconstruction versus Forecasting Utility}
\label{app:recutils}
Table~\ref{tab:app-recutils} gives the per-method reconstruction and forecasting numbers
behind the summary in \S\ref{sec:exp-opportunity}. Reconstruction error uses the artificially
hidden positions, which are known because the mask is synthetic; forecast gain is measured
against \textsc{Native Keep} on the same backbone. The discordant rate is the fraction of
method pairs whose reconstruction and forecasting rankings disagree.

\begin{table}[h]
  \caption{Full paired reconstruction-versus-utility table. The Spearman $\rho$ row compares
  the reconstruction ranking with the forecasting ranking within each backbone; a $\rho$
  close to zero means the two rankings carry different information.}
  \label{tab:app-recutils}
  \centering
  \small
  \setlength{\tabcolsep}{4pt}
\resizebox{\textwidth}{!}{\begin{tabular}{lccccc}
    \toprule
    Method & Rec.\ MSE $\downarrow$ & Rec.\ MAE $\downarrow$ & Forecast MASE $\downarrow$
           & Forecast Gain $\uparrow$ & Discordant Rate $\downarrow$ \\
    \midrule
    TOI                & \ph{REC_TOI_MSE}    & \ph{REC_TOI_MAE}    & \ph{REC_TOI_MASE}    & \ph{REC_TOI_GAIN}    & \multirow{6}{*}{\ph{DISCORDANT_RATE}} \\
    TOI-VSF            & \ph{REC_TOIVSF_MSE} & \ph{REC_TOIVSF_MAE} & \ph{REC_TOIVSF_MASE} & \ph{REC_TOIVSF_GAIN} & \\
    GIMCC              & \ph{REC_GIMCC_MSE}  & \ph{REC_GIMCC_MAE}  & \ph{REC_GIMCC_MASE}  & \ph{REC_GIMCC_GAIN}  & \\
    SRDI               & \ph{REC_SRDI_MSE}   & \ph{REC_SRDI_MAE}   & \ph{REC_SRDI_MASE}   & \ph{REC_SRDI_GAIN}   & \\
    ChannelTokenFormer & \ph{REC_CTF_MSE}    & \ph{REC_CTF_MAE}    & \ph{REC_CTF_MASE}    & \ph{REC_CTF_GAIN}    & \\
    \rowcolor{bestgray}\introact{} & \ph{REC_OURS_MSE} & \ph{REC_OURS_MAE} & \ph{REC_OURS_MASE} & \ph{REC_OURS_GAIN} & \\
    \midrule
    \multicolumn{6}{l}{\footnotesize Spearman $\rho$ between reconstruction rank and
    forecasting rank: Bolt \ph{RHO_BOLT}, TimesFM \ph{RHO_TF}, Chronos-2 \ph{RHO_CH2}} \\
    \bottomrule
  \end{tabular}}
\end{table}

\subsection{Governance Diagnostics}
\label{app:governance-fig}
Figure~\ref{fig:governance} visualises the diagnostics tabulated in
Table~\ref{tab:opportunity}. Panel~(b) is a sign-agreement check rather than a calibration
curve: the ranking score is used only to order actions, and we do not claim that its magnitude
is a calibrated probability.

\begin{figure}[h]
  \centering
  \includegraphics[width=0.84\textwidth]{figure/fig4_governance_diagnostics.pdf}
  \caption{\textbf{(a)}~Harmful loss by method. \textbf{(b)}~Realised utility against the
  binned ranking score with 95\% intervals. Panel~(b) is a sign agreement check, not a
  calibration claim.}
  \label{fig:governance}
\end{figure}


\subsection{Action-Opportunity Strata}
\label{app:opportunity}
The strata used in \S\ref{sec:exp-opportunity} are defined by thresholds on the oracle
opportunity $\Delta_i = L_{\textsc{Keep},i} - L_i^{\star}$ of \eqref{eq:opportunity}. The
thresholds are fitted on the gate split only and then applied unchanged to the evaluation
episodes, so no evaluation episode participates in choosing them. Table~\ref{tab:app-opp-sizes}
records the realised stratum sizes and the opportunity distribution, and
Table~\ref{tab:app-opp-sens} repeats the stratum-level MASE under alternative boundaries, so
that a reader can see how much of the effect depends on where the boundary was drawn.

\begin{table}[h]
  \caption{Action-opportunity strata on the evaluation episodes. Boundaries are the TRAIN-gate
  thresholds, applied unchanged. Counts are of legal parents, not of mask variants.}
  \label{tab:app-opp-sizes}
  \centering
  \small
  \setlength{\tabcolsep}{4pt}
\resizebox{\textwidth}{!}{\begin{tabular}{lcccc}
    \toprule
    Stratum & Parents & Share of episodes & Median $\Delta_i$ & 90th pct.\ $\Delta_i$ \\
    \midrule
    No-op             & \ph{STRATUM_NOOP_N}  & \ph{STRATUM_NOOP_SHARE}  & \ph{STRATUM_NOOP_MED}  & \ph{STRATUM_NOOP_P90} \\
    Low opportunity   & \ph{STRATUM_LOW_N}   & \ph{STRATUM_LOW_SHARE}   & \ph{STRATUM_LOW_MED}   & \ph{STRATUM_LOW_P90} \\
    High opportunity  & \ph{STRATUM_HIGH_N}  & \ph{STRATUM_HIGH_SHARE}  & \ph{STRATUM_HIGH_MED}  & \ph{STRATUM_HIGH_P90} \\
    \bottomrule
  \end{tabular}}
\end{table}

\begin{table}[h]
  \caption{Stratum-boundary sensitivity. The no-op tolerance and the low/high boundary are
  shifted by the indicated factor of the TRAIN-gate threshold; every other setting is
  unchanged. MASE is source-macro averaged within the resulting strata.}
  \label{tab:app-opp-sens}
  \centering
  \small
  \setlength{\tabcolsep}{4pt}
\resizebox{\textwidth}{!}{\begin{tabular}{lcccc}
    \toprule
    Boundary factor & No-op MASE $\downarrow$ & Low MASE $\downarrow$ & High MASE $\downarrow$ & Overall MASE $\downarrow$ \\
    \midrule
    $\times 0.5$ & \ph{OPPSENS_HALF_NOOP}  & \ph{OPPSENS_HALF_LOW}  & \ph{OPPSENS_HALF_HIGH}  & \ph{OPPSENS_HALF_OVR} \\
    $\times 1.0$ & \ph{OPPSENS_ONE_NOOP}   & \ph{OPPSENS_ONE_LOW}   & \ph{OPPSENS_ONE_HIGH}   & \ph{OPPSENS_ONE_OVR} \\
    $\times 2.0$ & \ph{OPPSENS_TWO_NOOP}   & \ph{OPPSENS_TWO_LOW}   & \ph{OPPSENS_TWO_HIGH}   & \ph{OPPSENS_TWO_OVR} \\
    \bottomrule
  \end{tabular}}
\end{table}

\begin{figure}[h]
  \centering
  \includegraphics[width=0.62\textwidth]{figure/fig3_action_opportunity.pdf}
  \caption{Action choice matters where the opportunity exists. The horizontal axis is the
  oracle opportunity $\Delta_i$ of an episode, binned from TRAIN-only thresholds; the
  vertical axis is the improvement of a method over \textsc{Keep}. The no-op stratum is where
  an intervention can only lose, and the high-opportunity stratum is where selection has
  something to gain. The panel is a placeholder until the main experiment produces the
  underlying CSV.}
  \label{fig:opportunity}
\end{figure}

\subsection{Extended Comparison}
\label{app:extended}
VIDA~\citep{liang2025vida} is a closely related variable-subset-forecasting framework that
reframes the problem as cross-domain knowledge transfer. It is not part of the five core
baselines: the core set is fixed at five published methods before any result is seen, and
VIDA's adaptation procedure targets a training regime that differs from our frozen-backbone
contract. We therefore report it as an extended comparison rather than substituting it for a
core baseline. Table~\ref{tab:app-vida} gives the comparison on the main protocol.

If a core baseline turns out not to admit a legal adaptation to our protocol, we report that
to the reader instead of replacing it silently, and any change to the core set is an explicit
protocol revision rather than an editorial one.

\begin{table}[h]
  \caption{Extended comparison on the main protocol ($10\%$ severity, source-macro MASE,
  lower is better). VIDA is reported for completeness and is excluded from the core average
  ranks of Table~\ref{tab:main}.}
  \label{tab:app-vida}
  \centering
  \small
  \setlength{\tabcolsep}{4pt}
  \begin{tabular}{lccccc}
    \toprule
    Method & Bolt $\downarrow$ & TimesFM $\downarrow$ & Chronos-2 $\downarrow$ & Overall $\downarrow$ & Rank $\downarrow$ \\
    \midrule
    VIDA           & \ph{VIDA_BOLT} & \ph{VIDA_TF} & \ph{VIDA_CH2} & \ph{VIDA_OVR} & \ph{VIDA_RANK} \\
    \introact{}    & \ph{OURS_BOLT} & \ph{OURS_TF} & \ph{OURS_CH2} & \ph{OURS_OVR} & \ph{OURS_RANK} \\
    \bottomrule
  \end{tabular}
\end{table}

\section{Robustness Breakdown}
\label{app:robustness}
This section expands \S\ref{sec:exp-robust}: the complete severity tables, the per-pattern
results, and the severity trend.

\subsection{Severity Tables}
\label{app:severity}
Tables~\ref{tab:app-sev10}--\ref{tab:app-sev50} give the complete severity sweep. No method
is retuned between severities: \introact{} reuses the $K$ and $\beta$ frozen after the
gate, and every baseline reuses the models and rules fixed on TRAIN. The $10\%$ column is
the main-text condition; the $30\%$ and $50\%$ columns exist to test whether the frozen
policy degrades faster or slower than the baselines, not to select a new configuration.

\begin{table}[h]
  \caption{Complete results at $10\%$ severity, all four metrics, source-macro averaged.
  Ranks are over deployable methods only and exclude the oracle row.}
  \label{tab:app-sev10}
  \centering
  \small
  \setlength{\tabcolsep}{4pt}
\resizebox{\textwidth}{!}{\begin{tabular}{lccccc}
    \toprule
    Method & MASE $\downarrow$ & MSE $\downarrow$ & MAE $\downarrow$ & RMSE $\downarrow$ & Rank $\downarrow$ \\
    \midrule
    \textsc{Native Keep} & \ph{SEV10-KEEP-MASE} & \ph{SEV10-KEEP-MSE} & \ph{SEV10-KEEP-MAE} & \ph{SEV10-KEEP-RMSE} & \ph{SEV10-KEEP-RANK} \\
    \textsc{Best Fixed}  & \ph{SEV10-BF-MASE}   & \ph{SEV10-BF-MSE}   & \ph{SEV10-BF-MAE}   & \ph{SEV10-BF-RMSE}   & \ph{SEV10-BF-RANK} \\
    R2-CART              & \ph{SEV10-R2-MASE}   & \ph{SEV10-R2-MSE}   & \ph{SEV10-R2-MAE}   & \ph{SEV10-R2-RMSE}   & \ph{SEV10-R2-RANK} \\
    TOI                   & \ph{SEV10_TOI_MASE}& \ph{SEV10_TOI_MSE}& \ph{SEV10_TOI_MAE}& \ph{SEV10_TOI_RMSE}& \ph{SEV10_TOI_RANK} \\
    TOI-VSF               & \ph{SEV10_TOIVSF_MASE} & \ph{SEV10_TOIVSF_MSE} & \ph{SEV10_TOIVSF_MAE} & \ph{SEV10_TOIVSF_RMSE} & \ph{SEV10_TOIVSF_RANK} \\
    GIMCC                 & \ph{SEV10_GIMCC_MASE}  & \ph{SEV10_GIMCC_MSE}  & \ph{SEV10_GIMCC_MAE}  & \ph{SEV10_GIMCC_RMSE}  & \ph{SEV10_GIMCC_RANK} \\
    SRDI                  & \ph{SEV10_SRDI_MASE}   & \ph{SEV10_SRDI_MSE}   & \ph{SEV10_SRDI_MAE}   & \ph{SEV10_SRDI_RMSE}   & \ph{SEV10_SRDI_RANK} \\
    ChannelTokenFormer    & \ph{SEV10_CTF_MASE} & \ph{SEV10_CTF_MSE} & \ph{SEV10_CTF_MAE} & \ph{SEV10_CTF_RMSE} & \ph{SEV10_CTF_RANK} \\
    \midrule
    \rowcolor{bestgray}\introact{} & \ph{SEV10-OURS-MASE} & \ph{SEV10-OURS-MSE} & \ph{SEV10-OURS-MAE} & \ph{SEV10-OURS-RMSE} & \ph{SEV10-OURS-RANK} \\
    \midrule
    \oracle{Catalog Oracle (diagnostic)} & \oracle{\ph{SEV10-ORACLE-MASE}} & \oracle{\ph{SEV10-ORACLE-MSE}} & \oracle{\ph{SEV10-ORACLE-MAE}} & \oracle{\ph{SEV10-ORACLE-RMSE}} & \oracle{---} \\
    \bottomrule
  \end{tabular}}
\end{table}

\begin{table}[h]
  \caption{Complete results at $30\%$ severity. Identical frozen configuration as the
  $10\%$ table; nothing is retuned.}
  \label{tab:app-sev30}
  \centering
  \small
  \setlength{\tabcolsep}{4pt}
\resizebox{\textwidth}{!}{\begin{tabular}{lccccc}
    \toprule
    Method & MASE $\downarrow$ & MSE $\downarrow$ & MAE $\downarrow$ & RMSE $\downarrow$ & Rank $\downarrow$ \\
    \midrule
    \textsc{Native Keep} & \ph{SEV30-KEEP-MASE} & \ph{SEV30-KEEP-MSE} & \ph{SEV30-KEEP-MAE} & \ph{SEV30-KEEP-RMSE} & \ph{SEV30-KEEP-RANK} \\
    \textsc{Best Fixed}  & \ph{SEV30-BF-MASE}   & \ph{SEV30-BF-MSE}   & \ph{SEV30-BF-MAE}   & \ph{SEV30-BF-RMSE}   & \ph{SEV30-BF-RANK} \\
    R2-CART              & \ph{SEV30-R2-MASE}   & \ph{SEV30-R2-MSE}   & \ph{SEV30-R2-MAE}   & \ph{SEV30-R2-RMSE}   & \ph{SEV30-R2-RANK} \\
    TOI                   & \ph{SEV30_TOI_MASE}& \ph{SEV30_TOI_MSE}& \ph{SEV30_TOI_MAE}& \ph{SEV30_TOI_RMSE}& \ph{SEV30_TOI_RANK} \\
    TOI-VSF               & \ph{SEV30_TOIVSF_MASE} & \ph{SEV30_TOIVSF_MSE} & \ph{SEV30_TOIVSF_MAE} & \ph{SEV30_TOIVSF_RMSE} & \ph{SEV30_TOIVSF_RANK} \\
    GIMCC                 & \ph{SEV30_GIMCC_MASE}  & \ph{SEV30_GIMCC_MSE}  & \ph{SEV30_GIMCC_MAE}  & \ph{SEV30_GIMCC_RMSE}  & \ph{SEV30_GIMCC_RANK} \\
    SRDI                  & \ph{SEV30_SRDI_MASE}   & \ph{SEV30_SRDI_MSE}   & \ph{SEV30_SRDI_MAE}   & \ph{SEV30_SRDI_RMSE}   & \ph{SEV30_SRDI_RANK} \\
    ChannelTokenFormer    & \ph{SEV30_CTF_MASE} & \ph{SEV30_CTF_MSE} & \ph{SEV30_CTF_MAE} & \ph{SEV30_CTF_RMSE} & \ph{SEV30_CTF_RANK} \\
    \midrule
    \rowcolor{bestgray}\introact{} & \ph{SEV30-OURS-MASE} & \ph{SEV30-OURS-MSE} & \ph{SEV30-OURS-MAE} & \ph{SEV30-OURS-RMSE} & \ph{SEV30-OURS-RANK} \\
    \bottomrule
  \end{tabular}}
\end{table}

\begin{table}[h]
  \caption{Complete results at $50\%$ severity. This condition is a stress test of the
  frozen policy, not a condition the method was tuned on.}
  \label{tab:app-sev50}
  \centering
  \small
  \setlength{\tabcolsep}{4pt}
\resizebox{\textwidth}{!}{\begin{tabular}{lccccc}
    \toprule
    Method & MASE $\downarrow$ & MSE $\downarrow$ & MAE $\downarrow$ & RMSE $\downarrow$ & Rank $\downarrow$ \\
    \midrule
    \textsc{Native Keep} & \ph{SEV50-KEEP-MASE} & \ph{SEV50-KEEP-MSE} & \ph{SEV50-KEEP-MAE} & \ph{SEV50-KEEP-RMSE} & \ph{SEV50-KEEP-RANK} \\
    \textsc{Best Fixed}  & \ph{SEV50-BF-MASE}   & \ph{SEV50-BF-MSE}   & \ph{SEV50-BF-MAE}   & \ph{SEV50-BF-RMSE}   & \ph{SEV50-BF-RANK} \\
    R2-CART              & \ph{SEV50-R2-MASE}   & \ph{SEV50-R2-MSE}   & \ph{SEV50-R2-MAE}   & \ph{SEV50-R2-RMSE}   & \ph{SEV50-R2-RANK} \\
    TOI                   & \ph{SEV50_TOI_MASE}& \ph{SEV50_TOI_MSE}& \ph{SEV50_TOI_MAE}& \ph{SEV50_TOI_RMSE}& \ph{SEV50_TOI_RANK} \\
    TOI-VSF               & \ph{SEV50_TOIVSF_MASE} & \ph{SEV50_TOIVSF_MSE} & \ph{SEV50_TOIVSF_MAE} & \ph{SEV50_TOIVSF_RMSE} & \ph{SEV50_TOIVSF_RANK} \\
    GIMCC                 & \ph{SEV50_GIMCC_MASE}  & \ph{SEV50_GIMCC_MSE}  & \ph{SEV50_GIMCC_MAE}  & \ph{SEV50_GIMCC_RMSE}  & \ph{SEV50_GIMCC_RANK} \\
    SRDI                  & \ph{SEV50_SRDI_MASE}   & \ph{SEV50_SRDI_MSE}   & \ph{SEV50_SRDI_MAE}   & \ph{SEV50_SRDI_RMSE}   & \ph{SEV50_SRDI_RANK} \\
    ChannelTokenFormer    & \ph{SEV50_CTF_MASE} & \ph{SEV50_CTF_MSE} & \ph{SEV50_CTF_MAE} & \ph{SEV50_CTF_RMSE} & \ph{SEV50_CTF_RANK} \\
    \midrule
    \rowcolor{bestgray}\introact{} & \ph{SEV50-OURS-MASE} & \ph{SEV50-OURS-MSE} & \ph{SEV50-OURS-MAE} & \ph{SEV50-OURS-RMSE} & \ph{SEV50-OURS-RANK} \\
    \bottomrule
  \end{tabular}}
\end{table}

\subsection{Per-Pattern Results}
\label{app:pattern}
Table~\ref{tab:app-pattern-res} separates the four patterns. Point missing is expected to
be the easiest regime and tail missing the hardest, because the missing block is adjacent
to the forecast origin and therefore removes exactly the observations a forecaster would
rely on most. Any claim about a pattern is a claim about this table, not about the pooled
main table.

\begin{table}[h]
  \caption{Per-pattern MASE at $10\%$ severity, source-macro averaged over the three
  backbones and both horizons. The last column is the average rank over the four patterns.}
  \label{tab:app-pattern-res}
  \centering
  \small
  \setlength{\tabcolsep}{4pt}
\resizebox{\textwidth}{!}{\begin{tabular}{lccccc}
    \toprule
    Method & P1 Point & P2 Target Block & P3 Shared Block & P4 Tail & Avg.\ Rank $\downarrow$ \\
    \midrule
    \textsc{Native Keep} & \ph{P1_KEEP} & \ph{P2_KEEP} & \ph{P3_KEEP} & \ph{P4_KEEP} & \ph{PR_KEEP} \\
    \textsc{Best Fixed}  & \ph{P1_BF}   & \ph{P2_BF}   & \ph{P3_BF}   & \ph{P4_BF}   & \ph{PR_BF} \\
    R2-CART              & \ph{P1_R2}   & \ph{P2_R2}   & \ph{P3_R2}   & \ph{P4_R2}   & \ph{PR_R2} \\
    TOI                   & \ph{P1_TOI}& \ph{P2_TOI}& \ph{P3_TOI}& \ph{P4_TOI}& \ph{PR_TOI} \\
    TOI-VSF               & \ph{P1_TOIVSF} & \ph{P2_TOIVSF} & \ph{P3_TOIVSF} & \ph{P4_TOIVSF} & \ph{PR_TOIVSF} \\
    GIMCC                 & \ph{P1_GIMCC}  & \ph{P2_GIMCC}  & \ph{P3_GIMCC}  & \ph{P4_GIMCC}  & \ph{PR_GIMCC} \\
    SRDI                  & \ph{P1_SRDI}   & \ph{P2_SRDI}   & \ph{P3_SRDI}   & \ph{P4_SRDI}   & \ph{PR_SRDI} \\
    ChannelTokenFormer    & \ph{P1_CTF} & \ph{P2_CTF} & \ph{P3_CTF} & \ph{P4_CTF} & \ph{PR_CTF} \\
    \midrule
    \rowcolor{bestgray}\introact{} & \ph{P1_OURS} & \ph{P2_OURS} & \ph{P3_OURS} & \ph{P4_OURS} & \ph{PR_OURS} \\
    \bottomrule
  \end{tabular}}
\end{table}

\subsection{Robustness Trend Figure}
\label{app:robust-fig}
Figure~\ref{fig:robust} shows the severity trend summarised by Table~\ref{tab:robust}. Each
panel is a different frozen backbone; the axes are identical across panels so that the
panels can be compared directly.

\begin{figure}[h]
  \centering
  \includegraphics[width=0.86\textwidth]{figure/fig5_robustness_missingness.pdf}
  \caption{Relative MASE improvement over \textsc{Native Keep} as missingness severity
  increases, one line per method, one panel per frozen backbone. All methods reuse the
  models and rules frozen on TRAIN, with no per-severity retuning.}
  \label{fig:robust}
\end{figure}


\section{Historical Support Results}
\label{app:support}
\label{app:support-details}
This section supports \S\ref{sec:exp-support}. Every method is re-fitted from scratch under
$25\%$, $50\%$, and $100\%$ of the replay-fit portion of TRAIN, using the identical split rule
and purge of Appendix~\ref{app:splits}. Table~\ref{tab:support} collects the compact reading
statement and Table~\ref{tab:app-support-full} reports all four metrics at each support level.
Table~\ref{tab:app-support-cost} separates the one-off offline fitting cost from the
per-request online cost at each level.

The bank-size check in Appendix~\ref{app:replaysize} is the corresponding internal ablation: it
subsamples the replay bank while holding the baseline training set fixed, whereas this section
restricts the available history for every method at the same time. The two are not
interchangeable, and a disagreement between them is informative rather than contradictory.

\begin{table}[h]
  \caption{Historical support efficiency. Every method is re-fitted under the same
  restriction on how much of replay-fit is available; MASE is source-macro averaged and lower
  is better. Offline cost is the one-off fitting cost of the method and is reported separately
  from the online cost of a request.}
  \label{tab:support}
  \centering
  \small
  \setlength{\tabcolsep}{3pt}
\resizebox{\textwidth}{!}{\begin{tabular}{lcccc}
    \toprule
    Method & 25\% support $\downarrow$ & 50\% support $\downarrow$ & 100\% support $\downarrow$ & Offline cost $\downarrow$ \\
    \midrule
    TOI                & \ph{SUPPORT25_TOI}    & \ph{SUPPORT50_TOI}    & \ph{SUPPORT100_TOI}    & \ph{OFFLINE_TOI} \\
    TOI-VSF            & \ph{SUPPORT25_TOIVSF} & \ph{SUPPORT50_TOIVSF} & \ph{SUPPORT100_TOIVSF} & \ph{OFFLINE_TOIVSF} \\
    GIMCC              & \ph{SUPPORT25_GIMCC}  & \ph{SUPPORT50_GIMCC}  & \ph{SUPPORT100_GIMCC}  & \ph{OFFLINE_GIMCC} \\
    SRDI               & \ph{SUPPORT25_SRDI}   & \ph{SUPPORT50_SRDI}   & \ph{SUPPORT100_SRDI}   & \ph{OFFLINE_SRDI} \\
    ChannelTokenFormer & \ph{SUPPORT25_CTF}    & \ph{SUPPORT50_CTF}    & \ph{SUPPORT100_CTF}    & \ph{OFFLINE_CTF} \\
    \midrule
    \rowcolor{bestgray}\introact{} & \ph{SUPPORT25_OURS} & \ph{SUPPORT50_OURS} & \ph{SUPPORT100_OURS} & \ph{OFFLINE_OURS} \\
    \bottomrule
  \end{tabular}}
\end{table}

\begin{table}[h]
  \caption{Historical support efficiency, all four metrics, source-macro averaged. Every
  method is re-fitted under the same restriction. Lower is better in every metric column.}
  \label{tab:app-support-full}
  \centering
  \small
  \setlength{\tabcolsep}{4pt}
\resizebox{\textwidth}{!}{\begin{tabular}{llcccc}
    \toprule
    Method & Support & MASE $\downarrow$ & MAE $\downarrow$ & RMSE $\downarrow$ & MSE $\downarrow$ \\
    \midrule
    TOI                & 25\% / 50\% / 100\% & \ph{TOI_S25_MASE} & \ph{TOI_S25_MAE} & \ph{TOI_S25_RMSE} & \ph{TOI_S25_MSE} \\
    TOI-VSF            & 25\% / 50\% / 100\% & \ph{TOIVSF_S25_MASE} & \ph{TOIVSF_S25_MAE} & \ph{TOIVSF_S25_RMSE} & \ph{TOIVSF_S25_MSE} \\
    GIMCC              & 25\% / 50\% / 100\% & \ph{GIMCC_S25_MASE} & \ph{GIMCC_S25_MAE} & \ph{GIMCC_S25_RMSE} & \ph{GIMCC_S25_MSE} \\
    SRDI               & 25\% / 50\% / 100\% & \ph{SRDI_S25_MASE} & \ph{SRDI_S25_MAE} & \ph{SRDI_S25_RMSE} & \ph{SRDI_S25_MSE} \\
    ChannelTokenFormer & 25\% / 50\% / 100\% & \ph{CTF_S25_MASE} & \ph{CTF_S25_MAE} & \ph{CTF_S25_RMSE} & \ph{CTF_S25_MSE} \\
    \introact{}        & 25\% / 50\% / 100\% & \ph{OURS_S25_MASE} & \ph{OURS_S25_MAE} & \ph{OURS_S25_RMSE} & \ph{OURS_S25_MSE} \\
    \bottomrule
  \end{tabular}}
\end{table}

\begin{table}[h]
  \caption{Offline and online cost by support level. Offline cost is the one-off fitting cost;
  online cost is per request and excludes cold start. The two are never summed into a single
  number, because they are paid on different schedules.}
  \label{tab:app-support-cost}
  \centering
  \small
  \setlength{\tabcolsep}{4pt}
\resizebox{\textwidth}{!}{\begin{tabular}{lcccccc}
    \toprule
    Method & Offline 25\% & Offline 50\% & Offline 100\% & Online 25\% & Online 50\% & Online 100\% \\
    \midrule
    TOI                & \ph{TOI_OFF25} & \ph{TOI_OFF50} & \ph{TOI_OFF100} & \ph{TOI_ON25} & \ph{TOI_ON50} & \ph{TOI_ON100} \\
    TOI-VSF            & \ph{TOIVSF_OFF25} & \ph{TOIVSF_OFF50} & \ph{TOIVSF_OFF100} & \ph{TOIVSF_ON25} & \ph{TOIVSF_ON50} & \ph{TOIVSF_ON100} \\
    GIMCC              & \ph{GIMCC_OFF25} & \ph{GIMCC_OFF50} & \ph{GIMCC_OFF100} & \ph{GIMCC_ON25} & \ph{GIMCC_ON50} & \ph{GIMCC_ON100} \\
    SRDI               & \ph{SRDI_OFF25} & \ph{SRDI_OFF50} & \ph{SRDI_OFF100} & \ph{SRDI_ON25} & \ph{SRDI_ON50} & \ph{SRDI_ON100} \\
    ChannelTokenFormer & \ph{CTF_OFF25} & \ph{CTF_OFF50} & \ph{CTF_OFF100} & \ph{CTF_ON25} & \ph{CTF_ON50} & \ph{CTF_ON100} \\
    \introact{}        & \ph{OURS_OFF25} & \ph{OURS_OFF50} & \ph{OURS_OFF100} & \ph{OURS_ON25} & \ph{OURS_ON50} & \ph{OURS_ON100} \\
    \bottomrule
  \end{tabular}}
\end{table}


\section{Ablation Details}
\label{app:ablation-details}
This section supports \S\ref{sec:exp-ablation}. Table~\ref{tab:ablation} reports the six
variants on the unified main protocol; Table~\ref{tab:app-ablation-full} reports the full
outcome vector; and Table~\ref{tab:app-ablation-stab} repeats the ablation under a different
deterministic ordering of the replay records, so that a reader can see whether a variant's
position is an artefact of the order in which records were consumed. The variant definitions
follow \S\ref{sec:exp-ablation} exactly.

\begin{table}[h]
  \caption{Ablation on the main protocol; $\checkmark$ marks a present component and
  ``partial'' a state group reduced to a subset. \emph{Replay sup.} is historical replay
  supervision, \emph{Regret rank} the regret-aware ranking objective, \emph{Pair w.} the
  pairwise comparison weights, \emph{Ep./Int.\ state} the two state groups, and
  \emph{Keep sel.} the availability of \textsc{Keep} in the decision space. Harmful loss is
  measured relative to \textsc{Native Keep}.}
  \label{tab:ablation}
  \centering
  \small
  \setlength{\tabcolsep}{3pt}
\resizebox{\textwidth}{!}{\begin{tabular}{lcccccccccc}
    \toprule
    & \multicolumn{5}{c}{Components} & \multicolumn{5}{c}{Outcome} \\
    \cmidrule(lr){2-6}\cmidrule(lr){7-11}
    Variant & Replay sup. & Regret rank & Pair w. & Ep./Int.\ state & Keep sel.
            & MASE $\downarrow$ & MAE $\downarrow$ & RMSE $\downarrow$ & Harmful loss $\downarrow$ & Rank $\downarrow$ \\
    \midrule
    Full \introact{}          & $\checkmark$ & $\checkmark$ & $\checkmark$ & $\checkmark$ & $\checkmark$ & \ph{FULL_MASE} & \ph{FULL_MAE} & \ph{FULL_RMSE} & \ph{FULL_HL} & \ph{FULL_RANK} \\
    A1 Gain regression        & $\checkmark$ &              & ---          & $\checkmark$ & $\checkmark$ & \ph{A1_MASE}   & \ph{A1_MAE}   & \ph{A1_RMSE}   & \ph{A1_HL}   & \ph{A1_RANK} \\
    A2 Unweighted ranking     & $\checkmark$ & $\checkmark$ &              & $\checkmark$ & $\checkmark$ & \ph{A2_MASE}   & \ph{A2_MAE}   & \ph{A2_RMSE}   & \ph{A2_HL}   & \ph{A2_RANK} \\
    A3 w/o intervention state & $\checkmark$ & $\checkmark$ & $\checkmark$ & partial      & $\checkmark$ & \ph{A3_MASE}   & \ph{A3_MAE}   & \ph{A3_RMSE}   & \ph{A3_HL}   & \ph{A3_RANK} \\
    A4 w/o context state      & $\checkmark$ & $\checkmark$ & $\checkmark$ & partial      & $\checkmark$ & \ph{A4_MASE}   & \ph{A4_MAE}   & \ph{A4_RMSE}   & \ph{A4_HL}   & \ph{A4_RANK} \\
    A5 Forced intervention    & $\checkmark$ & $\checkmark$ & $\checkmark$ & $\checkmark$ &              & \ph{A5_MASE}   & \ph{A5_MAE}   & \ph{A5_RMSE}   & \ph{A5_HL}   & \ph{A5_RANK} \\
    \bottomrule
  \end{tabular}}
\end{table}

\begin{table}[h]
  \caption{Ablation, full outcome vector. Harmful loss is measured relative to
  \textsc{Native Keep} and is reported both unconditionally and conditional on an intervention
  having been executed. Calls is the mean number of frozen-model forecasts per request.}
  \label{tab:app-ablation-full}
  \centering
  \small
  \setlength{\tabcolsep}{4pt}
\resizebox{\textwidth}{!}{\begin{tabular}{lccccccc}
    \toprule
    Variant & MASE $\downarrow$ & MAE $\downarrow$ & RMSE $\downarrow$ & MSE $\downarrow$ & Harmful loss $\downarrow$ & Conditional HIR $\downarrow$ & Calls $\downarrow$ \\
    \midrule
    Full \introact{}          & \ph{AF_FULL_MASE} & \ph{AF_FULL_MAE} & \ph{AF_FULL_RMSE} & \ph{AF_FULL_MSE} & \ph{AF_FULL_HL} & \ph{AF_FULL_CHIR} & \ph{AF_FULL_CALLS} \\
    A1 Gain regression        & \ph{AF_A1_MASE}   & \ph{AF_A1_MAE}   & \ph{AF_A1_RMSE}   & \ph{AF_A1_MSE}   & \ph{AF_A1_HL}   & \ph{AF_A1_CHIR}   & \ph{AF_A1_CALLS} \\
    A2 Unweighted ranking     & \ph{AF_A2_MASE}   & \ph{AF_A2_MAE}   & \ph{AF_A2_RMSE}   & \ph{AF_A2_MSE}   & \ph{AF_A2_HL}   & \ph{AF_A2_CHIR}   & \ph{AF_A2_CALLS} \\
    A3 w/o intervention state & \ph{AF_A3_MASE}   & \ph{AF_A3_MAE}   & \ph{AF_A3_RMSE}   & \ph{AF_A3_MSE}   & \ph{AF_A3_HL}   & \ph{AF_A3_CHIR}   & \ph{AF_A3_CALLS} \\
    A4 w/o context state      & \ph{AF_A4_MASE}   & \ph{AF_A4_MAE}   & \ph{AF_A4_RMSE}   & \ph{AF_A4_MSE}   & \ph{AF_A4_HL}   & \ph{AF_A4_CHIR}   & \ph{AF_A4_CALLS} \\
    A5 Forced intervention    & \ph{AF_A5_MASE}   & \ph{AF_A5_MAE}   & \ph{AF_A5_RMSE}   & \ph{AF_A5_MSE}   & \ph{AF_A5_HL}   & \ph{AF_A5_CHIR}   & \ph{AF_A5_CALLS} \\
    \bottomrule
  \end{tabular}}
\end{table}

\begin{table}[h]
  \caption{Ordering stability of the ablation. The whole ablation is repeated after a
  deterministic reshuffle of the replay records; the column reports the Spearman correlation
  between the two rankings of the six variants, per backbone. A high value means the variant
  ordering is not an artefact of record order.}
  \label{tab:app-ablation-stab}
  \centering
  \small
  \setlength{\tabcolsep}{4pt}
  \begin{tabular}{lcc}
    \toprule
    Backbone & Spearman $\rho$ between orderings $\uparrow$ & Rank of Full variant \\
    \midrule
    Bolt      & \ph{ABL_STAB_BOLT} & \ph{ABL_RANKFULL_BOLT} \\
    TimesFM   & \ph{ABL_STAB_TF}   & \ph{ABL_RANKFULL_TF} \\
    Chronos-2 & \ph{ABL_STAB_CH2}  & \ph{ABL_RANKFULL_CH2} \\
    \bottomrule
  \end{tabular}
\end{table}

\section{Runtime and Failure Audit}
\label{app:calls}
Table~\ref{tab:app-efficiency} gives the full cost breakdown behind the summary in
\S\ref{sec:exp-main}.

\begin{table}[t]
  \caption{Deployment cost for all methods, offline and online kept separate. Latency percentiles are
  per request and exclude cold start, reported on its own line. Lower is better except for
  the call count.}
  \label{tab:app-efficiency}
  \centering
  \small
  \setlength{\tabcolsep}{4pt}
\resizebox{\textwidth}{!}{\begin{tabular}{lcccccc}
    \toprule
    Method & Offline train/search & Candidate constr. & TSFM calls/req.\ & Mean lat.\ $\downarrow$ & P95 $\downarrow$ & Max $\downarrow$ \\
    \midrule
    TOI                & \ph{TOI_OFF}    & \ph{TOI_CAND}    & \ph{TOI_CALLS}    & \ph{TOI_MEAN}    & \ph{TOI_P95}    & \ph{TOI_MAX} \\
    TOI-VSF            & \ph{TOIVSF_OFF} & \ph{TOIVSF_CAND} & \ph{TOIVSF_CALLS} & \ph{TOIVSF_MEAN} & \ph{TOIVSF_P95} & \ph{TOIVSF_MAX} \\
    GIMCC              & \ph{GIMCC_OFF}  & \ph{GIMCC_CAND}  & \ph{GIMCC_CALLS}  & \ph{GIMCC_MEAN}  & \ph{GIMCC_P95}  & \ph{GIMCC_MAX} \\
    SRDI               & \ph{SRDI_OFF}   & \ph{SRDI_CAND}   & \ph{SRDI_CALLS}   & \ph{SRDI_MEAN}   & \ph{SRDI_P95}   & \ph{SRDI_MAX} \\
    ChannelTokenFormer & \ph{CTF_OFF}    & \ph{CTF_CAND}    & \ph{CTF_CALLS}    & \ph{CTF_MEAN}    & \ph{CTF_P95}    & \ph{CTF_MAX} \\
    \rowcolor{bestgray}\introact{} & \ph{OURS_OFF} & \ph{OURS_CAND} & \ph{OURS_CALLS} & \ph{OURS_MEAN} & \ph{OURS_P95} & \ph{OURS_MAX} \\
    \midrule
    \multicolumn{7}{l}{\footnotesize Cold start: \ph{COLDSTART} \quad
    Failure/timeout share: \ph{FAILRATE}} \\
    \bottomrule
  \end{tabular}}
\end{table}


Table~\ref{tab:app-calls} separates offline from online cost and separates successes from
failures. Failures and timeouts are never counted as successes and are never dropped from
the denominator; a request that fails to produce a forecast is reported as a failure, not as
a zero-cost success. Under the frozen state definition of \S\ref{sec:method-scoring} the
scorer reads no forecast, so a request costs exactly one frozen-model call: the call is made
for the action the ranking selects, and when that action is \textsc{Keep} the same call
produces the returned reference forecast. Candidate materialisation and scoring are
forecast-free and are booked separately.

\begin{table}[h]
  \caption{Frozen-model call accounting. Latency is reported per request and excludes cold
  start, which is reported on its own line; the offline replay-bank build and the one-off
  selector fit are amortised and reported separately. Candidate materialisation and scoring
  are forecast-free, so they do not appear in the call column.}
  \label{tab:app-calls}
  \centering
  \small
  \setlength{\tabcolsep}{4pt}
\resizebox{\textwidth}{!}{\begin{tabular}{lccc}
    \toprule
    Outcome & Requests & Share & Frozen-model calls per request \\
    \midrule
    \textsc{Keep} selected (input untouched)   & \ph{CALLS_KEEP_N}    & \ph{CALLS_KEEP_SHARE}    & 1 \\
    Intervention selected (input changed)      & \ph{CALLS_ACT_N}     & \ph{CALLS_ACT_SHARE}     & 1 \\
    Failed before any forecast                 & \ph{CALLS_FAIL_N}    & \ph{CALLS_FAIL_SHARE}    & \ph{CALLS_FAIL_COST} \\
    Timed out after the forecast               & \ph{CALLS_TIMEOUT_N} & \ph{CALLS_TIMEOUT_SHARE} & \ph{CALLS_TIMEOUT_COST} \\
    \midrule
    Offline replay-bank build (amortised)      & --- & --- & \ph{CALLS_OFFLINE} \\
    Offline selector fit (amortised)           & --- & --- & \ph{CALLS_SELECTOR} \\
    Cold start (first request of a process)    & \ph{CALLS_COLD_N} & --- & \ph{CALLS_COLD_LATENCY} \\
    \bottomrule
  \end{tabular}}
\end{table}


\section{Development Negative Results}
\label{app:development}
This section records what was tried in the development of this line of work and rejected, and
why. It exists because the final formulation is easier to read against the alternatives it
replaced, and because omitting it would misrepresent how the method was reached. Nothing in
this section is an independent confirmation of any method, and no number here is comparable
with any table in the main text.

\subsection{Local Gain Averaging with Lower-Confidence Gating}
\label{app:r5}
An earlier configuration in this line of work attempted to improve intervention ranking by
projecting estimated gain vectors onto a convex feasible set, under a Frank--Wolfe
solver~\citep{boyd2004convex,jaggi2013revisiting}. The projection has a clean property ---
for an exact projection $z$ of a score vector $r$, one has
$\lVert z - g \rVert^2 \le \lVert r - g \rVert^2 - \lVert r - z \rVert^2$ for any feasible
true gain vector $g$ --- but that property controls the squared error of the whole vector
only. It does not guarantee per-coordinate improvement, does not preserve action ranking,
and does not imply a better MASE.

In that cycle the projection did reduce the squared error of the estimated gain vector
while failing to improve, and in one backbone degrading, the final forecasting metric.
Those measurements are recorded here because they motivate the design decision in
\S\ref{sec:method-scoring}: continuous score calibration is insufficient for discrete
intervention selection, which is why \introact{} scores actions with an explicitly
conservative mean-minus-uncertainty quantity and abstains rather than refining a score.
The recorded measurements belong to the preceding protocol (a small development split of
26 parents with 156 correlated variants, two backbones, and a different action budget
accounting). They are reproduced below verbatim from the frozen ledger and must not be
compared numerically with any table in this paper.

\begin{table}[h]
  \caption{Preceding-cycle measurements, retained as a negative result. Different protocol
  from the present paper: 26 development parents, 156 correlated variants, different
  budget accounting. Values are quoted from the frozen ledger and are \emph{not}
  comparable with Tables~\ref{tab:main}--\ref{tab:ablation}.}
  \label{tab:r5negative}
  \centering
  \small
  \begin{tabular}{llcc}
    \toprule
    Backbone & Configuration & MASE $\downarrow$ & Note \\
    \midrule
    Bolt     & Native KEEP (no governance)   & 1.258454 & reference \\
    Bolt     & Free reference policy         & 1.152415 & equals the full candidate \\
    Bolt     & Projection candidate          & 1.152415 & no gain over the free reference \\
    TimesFM  & Native KEEP (no governance)   & 1.139837 & reference \\
    TimesFM  & Strong simple TRAIN control   & 1.069398 & strongest simple rule \\
    TimesFM  & Projection candidate          & 1.074244 & worse than the simple control \\
    \bottomrule
  \end{tabular}
\end{table}

Two consequences of that cycle are load-bearing for the present paper. First, the strongest
simple selector available at the time (\textbf{R2-CART} in the present terminology) was
competitive with, and on one backbone better than, the more elaborate candidate; that is
why R2-CART is carried forward as a mandatory control in every main-text table. Second,
reducing estimation error on a continuous score was not sufficient to improve discrete
selection; that is why the lower-confidence gate in \S\ref{sec:method-scoring} is evaluated
directly on harmful intervention rate and harmful loss, not only on MASE.

\subsection{Sensitivity of the Development-Cycle Hyperparameters}
\label{app:ksens}
The rejected formulation of Appendix~\ref{app:development} had exactly two tuned quantities:
a neighbourhood size $K \in \{8, 16, 32\}$ and a conservativeness level
$\beta \in \{0, 1, 1.64\}$. Table~\ref{tab:app-ksens} reports the full $3 \times 3$ grid on the
gate split, so that a reader can see how the rejected configuration was chosen and how flat
the surface was. The grid is retained as development evidence; it does not describe the frozen
configuration of this paper, whose hyperparameters are frozen on the same gate split and
reported in \S\ref{sec:exp-main}.

\begin{table}[h]
  \caption{Sensitivity of the rejected development-cycle configuration on the gate split.
  Each cell reports MASE and the harmful intervention rate. The cell that was carried into
  that cycle's configuration is marked with a dagger. These cells belong to the rejected
  formulation and are \emph{not} comparable with any table in the main text.}
  \label{tab:app-ksens}
  \centering
  \small
  \setlength{\tabcolsep}{4pt}
  \begin{tabular}{lccc}
    \toprule
    & $\beta = 0$ & $\beta = 1$ & $\beta = 1.64$ \\
    \midrule
    $K = 8$  & \ph{K8_B0}  & \ph{K8_B1}  & \ph{K8_B164} \\
    $K = 16$ & \ph{K16_B0} & \ph{K16_B1} & \ph{K16_B164} \\
    $K = 32$ & \ph{K32_B0} & \ph{K32_B1} & \ph{K32_B164} \\
    \bottomrule
  \end{tabular}
\end{table}


\subsection{The Sealed v4.4 TRAIN-Eval Batch}
\label{app:train-eval}
The preceding v4.4 cycle was run to completion on the TRAIN-only splits with the confirmation
set left sealed, in order to decide whether local gain averaging with a lower-confidence gate
should be promoted to the main experiment. It was not promoted. We report the batch because
the decision it motivated is part of the final method's justification.

Table~\ref{tab:app-train-eval} gives the source-macro MASE of the frozen v4.4 configuration and
of its own ablations, on the two development backbones. Four observations are load-bearing.
\textbf{(i)}~Counterfactual replay does carry downstream decision signal: the full
configuration is below \textsc{Native Keep} on both backbones. \textbf{(ii)}~The
lower-confidence gate lowers the harmful-intervention rate but pays for it in average accuracy:
removing it (\textsc{A4}) improves MASE on both backbones, so the gate is a governance
trade-off rather than an accuracy gain. \textbf{(iii)}~The local-matching component does not
carry independent evidence: the parametric gain regressors are at least as good as, and
usually better than, the full configuration. \textbf{(iv)}~Gain-estimation accuracy and action
ordering are not the same objective, which is why the final formulation optimises the ordering
directly rather than the point-wise error of a gain vector.

The batch is development evidence on TRAIN-only splits with the confirmation set unread. It is
not a promotion, it is not independent confirmation, and its numbers are not comparable with
Tables~\ref{tab:main}--\ref{tab:ablation}.

\begin{table}[h]
  \caption{Sealed v4.4 TRAIN-eval batch, retained as development evidence. Values are
  source-macro MASE from the development protocol; the confirmation set was not read. HIR is
  the harmful-intervention rate relative to \textsc{Native Keep}. These numbers belong to a
  different configuration and protocol and are \emph{not} comparable with any table in the main
  text.}
  \label{tab:app-train-eval}
  \centering
  \small
  \setlength{\tabcolsep}{4pt}
  \begin{tabular}{llcc}
    \toprule
    Backbone & Configuration & MASE $\downarrow$ & HIR $\downarrow$ \\
    \midrule
    Bolt     & Native KEEP            & 1.179735 & 0.000000 \\
    Bolt     & Best Fixed             & 1.179735 & 0.000000 \\
    Bolt     & R2-CART                & 1.178209 & 0.048913 \\
    Bolt     & Full v4.4              & 1.165800 & 0.084239 \\
    Bolt     & A4 (no gate)           & 1.142909 & 0.274457 \\
    Bolt     & A5 gain regression (ridge) & 1.140129 & 0.290761 \\
    Bolt     & A5 gain regression (CART)  & 1.153831 & 0.296196 \\
    \midrule
    TimesFM  & Native KEEP            & 1.242208 & 0.000000 \\
    TimesFM  & Best Fixed             & 1.103034 & 0.282609 \\
    TimesFM  & R2-CART                & 1.146845 & 0.078804 \\
    TimesFM  & Full v4.4              & 1.101261 & 0.187500 \\
    TimesFM  & A4 (no gate)           & 1.085381 & 0.339674 \\
    TimesFM  & A5 gain regression (ridge) & 1.088242 & 0.290761 \\
    TimesFM  & A5 gain regression (CART)  & 1.080967 & 0.334239 \\
    \bottomrule
  \end{tabular}
\end{table}

\section{Complete-Input Contract}
\label{app:complete}
Under a complete and valid input, \introact{} must return \textsc{Keep}. No forecasting
improvement is claimed in this regime; the purpose of this appendix is to document the
governance contract rather than to report a result. Table~\ref{tab:app-complete} lists the
regimes and the required behaviour in each.

\begin{table}[h]
  \caption{Governance contract by input regime. The complete-input rows are contract
  statements, not accuracy results: the expected behaviour is a no-op, and a non-no-op
  there would be a defect rather than an improvement.}
  \label{tab:app-complete}
  \centering
  \small
  \setlength{\tabcolsep}{4pt}
\resizebox{\textwidth}{!}{\begin{tabular}{lll}
    \toprule
    Input regime & Required behaviour & Claim \\
    \midrule
    Complete and valid input            & return \textsc{Keep} and the reference forecast & no-op contract, no accuracy claim \\
    Complete input, malformed request   & reject with an explicit status & no claim \\
    Unsupported action requested        & reject the action, keep the catalog closed & no claim \\
    Incomplete input, no strong evidence& abstain, return the reference forecast & abstention is a valid outcome \\
    Incomplete input, strong evidence   & execute exactly one action and return its forecast & evaluated in \S\ref{sec:exp-main} \\
    \bottomrule
  \end{tabular}}
\end{table}

\section{Financial Case Study}
\label{app:finance}
The two finance parents currently available are reported as a case study only. We do not
claim financial generalisation. Results are reported separately for the complete-observed
condition and for controlled deletion, and verified natural missingness is reported only
if it exists; where it does not, that is stated rather than substituted with synthetic
gaps. QLIKE is reported here and nowhere else. The finance series contain real jumps, such
as market closures; these are not smoothed and are not filled as if they were missing
observations.

\begin{table}[h]
  \caption{Finance case study. The complete-observed rows are the no-op check; the
  controlled-deletion rows carry the only quantitative comparison. Natural missingness is
  listed only where it has been verified.}
  \label{tab:app-finance}
  \centering
  \small
  \setlength{\tabcolsep}{4pt}
\resizebox{\textwidth}{!}{\begin{tabular}{llcccc}
    \toprule
    Parent & Condition & MASE $\downarrow$ & QLIKE $\downarrow$ & Harmful Rate & Note \\
    \midrule
    \ph{FIN_P1} & complete observed     & \ph{FIN-P1-COMP-MASE} & \ph{FIN-P1-COMP-QLIKE} & --- & no-op check \\
    \ph{FIN_P1} & controlled deletion   & \ph{FIN-P1-DEL-MASE}  & \ph{FIN-P1-DEL-QLIKE}  & \ph{FIN-P1-HIR} & \ph{FIN-P1-NOTE} \\
    \ph{FIN_P2} & complete observed     & \ph{FIN-P2-COMP-MASE} & \ph{FIN-P2-COMP-QLIKE} & --- & no-op check \\
    \ph{FIN_P2} & controlled deletion   & \ph{FIN-P2-DEL-MASE}  & \ph{FIN-P2-DEL-QLIKE}  & \ph{FIN-P2-HIR} & \ph{FIN-P2-NOTE} \\
    \midrule
    \multicolumn{6}{l}{\footnotesize Verified natural missingness: \ph{FIN-NATURAL-STATUS}} \\
    \bottomrule
  \end{tabular}}
\end{table}

\section{Additional Seeds and Sensitivity}
\label{app:seeds}
The main experiment uses one deterministic mask per parent. This appendix reruns a
pre-registered representative protocol with ten additional mask seeds to show that the
conclusions are stable with respect to mask realisation. The ten seeds are \emph{not}
treated as ten independent parents and are never used for significance testing: the
statistical unit remains the parent, and the seeds are read as a stability check on the
same parent set.

\begin{table}[h]
  \caption{Mask-seed stability. Each row is one additional mask realisation on the
  pre-registered representative parent subset. Spread across seeds is a stability
  diagnostic, not an additional sample of parents.}
  \label{tab:app-seeds}
  \centering
  \small
  \setlength{\tabcolsep}{4pt}
  \begin{tabular}{lccc}
    \toprule
    Mask seed & MASE $\downarrow$ & $\Delta$ vs \textsc{Keep} $\uparrow$ & Intervention Rate \\
    \midrule
    \ph{SEED_1}  & \ph{SEED1_MASE}  & \ph{SEED1_DELTA}  & \ph{SEED1_IR} \\
    \ph{SEED_2}  & \ph{SEED2_MASE}  & \ph{SEED2_DELTA}  & \ph{SEED2_IR} \\
    \ph{SEED_3}  & \ph{SEED3_MASE}  & \ph{SEED3_DELTA}  & \ph{SEED3_IR} \\
    \ph{SEED_4}  & \ph{SEED4_MASE}  & \ph{SEED4_DELTA}  & \ph{SEED4_IR} \\
    \ph{SEED_5}  & \ph{SEED5_MASE}  & \ph{SEED5_DELTA}  & \ph{SEED5_IR} \\
    \ph{SEED_6}  & \ph{SEED6_MASE}  & \ph{SEED6_DELTA}  & \ph{SEED6_IR} \\
    \ph{SEED_7}  & \ph{SEED7_MASE}  & \ph{SEED7_DELTA}  & \ph{SEED7_IR} \\
    \ph{SEED_8}  & \ph{SEED8_MASE}  & \ph{SEED8_DELTA}  & \ph{SEED8_IR} \\
    \ph{SEED_9}  & \ph{SEED9_MASE}  & \ph{SEED9_DELTA}  & \ph{SEED9_IR} \\
    \ph{SEED_10} & \ph{SEED10_MASE} & \ph{SEED10_DELTA} & \ph{SEED10_IR} \\
    \midrule
    Spread (max $-$ min) & \ph{SEED_SPREAD} & --- & \ph{SEED-IR-SPREAD} \\
    \bottomrule
  \end{tabular}
\end{table}
\subsection{Additional Horizons}
\label{app:horizons}
Longer horizons are reported here and are excluded from the main text, so that the main
comparison stays on the two horizons that every baseline can serve. These horizons are
reported for completeness; a conclusion that only holds at a longer horizon is not
presented as a main claim.

\begin{table}[h]
  \caption{Additional horizons beyond the two main ones, source-macro averaged, $10\%$
  severity. Reported for completeness and excluded from all main-text ranks.}
  \label{tab:app-horizons}
  \centering
  \small
  \setlength{\tabcolsep}{4pt}
\resizebox{\textwidth}{!}{\begin{tabular}{lccccc}
    \toprule
    Horizon & Method & MASE $\downarrow$ & MSE $\downarrow$ & MAE $\downarrow$ & RMSE $\downarrow$ \\
    \midrule
    \ph{HX1} & \introact{} & \ph{HX1-OURS-MASE} & \ph{HX1-OURS-MSE} & \ph{HX1-OURS-MAE} & \ph{HX1-OURS-RMSE} \\
    \ph{HX1} & \textsc{Native Keep} & \ph{HX1-KEEP-MASE} & \ph{HX1-KEEP-MSE} & \ph{HX1-KEEP-MAE} & \ph{HX1-KEEP-RMSE} \\
    \ph{HX2} & \introact{} & \ph{HX2-OURS-MASE} & \ph{HX2-OURS-MSE} & \ph{HX2-OURS-MAE} & \ph{HX2-OURS-RMSE} \\
    \ph{HX2} & \textsc{Native Keep} & \ph{HX2-KEEP-MASE} & \ph{HX2-KEEP-MSE} & \ph{HX2-KEEP-MAE} & \ph{HX2-KEEP-RMSE} \\
    \bottomrule
  \end{tabular}}
\end{table}
\subsection{Reproducibility Details}
\label{app:reprodetails}
Each record in the result files contains: source, parent, origin, horizon, pattern,
severity, mask hash, backbone, backbone revision, method, action, input hash, prediction
hash, forecast, target reference id, MASE, MSE, MAE, RMSE, reconstruction MSE where
applicable, runtime components, failure, fallback, code SHA, and resolved configuration.
Results are written to the logical result groups listed in Table~\ref{tab:app-groups}, each
in CSV, JSON, and Markdown form, with raw predictions retained alongside the aggregates.

\begin{table}[h]
  \caption{Result groups written by the pipeline. Every number in the main text and in the
  appendices is generated from one of these groups; nothing is transcribed by hand.}
  \label{tab:app-groups}
  \centering
  \small
  \setlength{\tabcolsep}{4pt}
\resizebox{\textwidth}{!}{\begin{tabular}{lll}
    \toprule
    Result group & Contents & Consumed by \\
    \midrule
    \texttt{v44\_protocol}       & frozen protocol, source manifests, split bounds & Appendix~\ref{app:replay} \\
    \texttt{v44\_replay\_bank}   & replay records $(z_i, a, g_{i,a})$ with input and prediction hashes & \S\ref{sec:method-replay} \\
    \texttt{v44\_baselines}      & baseline training runs, resolved hyperparameters, checkpoints & \S\ref{sec:exp-main} \\
    \texttt{v44\_train\_eval}    & TRAIN-eval acceptance check before freezing & Appendix~\ref{app:splits} \\
    \texttt{v44\_main\_results}  & main matrix rows, all methods, all cells & Tables~\ref{tab:main}, \ref{tab:app-src-bolt}--\ref{tab:app-src-ch2} \\
    \texttt{v44\_ablation}       & the six ablation variants & Table~\ref{tab:ablation} \\
    \texttt{v44\_robustness}     & severity, pattern, and mask-seed sweeps & Tables~\ref{tab:robust}, \ref{tab:app-sev10}--\ref{tab:app-sev50}, \ref{tab:app-seeds} \\
    \texttt{v44\_cost\_audit}    & call counts, latency, failures, timeouts & Table~\ref{tab:app-efficiency}, Appendix~\ref{app:calls} \\
    \bottomrule
  \end{tabular}}
\end{table}

The protocol is frozen before the main experiment and identified by commit hash. Any
change to the action catalog, the missingness protocol, the mask derivation rule, the split
boundaries, the state feature list, or the failure policy invalidates the freeze and
requires a new configuration identifier. Calibration and test labels are not read during
development, and no result in this paper is selected on them.

\end{document}
